\documentclass{article}
\usepackage[12pt]{extsizes}
\usepackage[a4paper, margin=0.5cm]{geometry}
\usepackage{indentfirst}
\usepackage[T1, T2A]{fontenc}
\usepackage[english, russian]{babel}
\usepackage{amsmath, amssymb, fleqn, cancel, graphicx, wrapfig, tabbing, color, cases}

\usepackage{amstext} % for \text macro
\usepackage{array}   % for \newcolumntype macro
\newcolumntype{L}{>{$}l<{$}} % math-mode version of "l" column type

\DeclareRobustCommand{\divby}{\mathrel{\vbox{\baselineskip.65ex\lineskiplimit0pt\hbox{.}\hbox{.}\hbox{.}}}}

\newenvironment{rcases}
    {\left.\begin{aligned}}
    {\end{aligned}\right\rbrace}

\usepackage{hyperref}
\hypersetup{
    colorlinks,
    citecolor=black,
    filecolor=black,
    linkcolor=black,
    urlcolor=black
}



\begin{document}

\begin{titlepage}
    \title{Дискретная математика}
    \author{Бабицкий Иван, М8О-111БВ-24}
    \date{2024/2025 учебный год, 1-й семестр}
    \maketitle
    \begin{center}
        
\includegraphics[scale=0.2]{images/mai.png}
    \end{center}
\end{titlepage}

\tableofcontents




\newpage
\paragraph{2024.09.17 (3, 4)}
\part{Теория множеств}

\section{Введение}

\subsection{Определение | Форма от $x$}
$P(x)$ ("форма от $x$") --- \{И, Л\} \\
$A =$ \{$x$ | $P(x) =$ И\} = \{$x$ | $P(x)$\} ("множество $A$ --- это множество таких $x$, что форма от $x$ истинна")

\subsubsection{Пример}
\begin{gather*}
    x^2 - 1 = 0 \\
    A = \{1, -1\} = \{x~|~x^2 - 1 = 0\}
\end{gather*}


\subsection{Свойства}
\begin{enumerate}
    \item $\{1, 2\} = \{2, 1, 1, 2, 1\}$ \\
    \item $A \subseteq B \Rightarrow \forall x \in A$ \\
    \item $A \subseteq B, B \subseteq C \Rightarrow A \subseteq C$ \\
    \item $A = B \Leftrightarrow A \subseteq B, B \subseteq A$
\end{enumerate}

\subsubsection{Пример}
\[
    \mathbb{N} \subseteq \mathbb{Z} \subseteq \mathbb{Q} \subseteq \mathbb{R} \subseteq \mathbb{C}
\]

\subsubsection{Пример}
\begin{gather*}
    A = \{1, 2\}, B = \{ \{1, 2\}, 3 \}, C = \{1, 2, 3\} \\
    1 \in A, 1 \notin B, A \subseteq C, B \neq C, A \in B
\end{gather*}


\subsection{Определение | Количество элементов в множестве}
\[
    m(A) = |A| = \#A
\]

\subsubsection{Пример}
\[
    m(\varnothing) = 0
\]


\subsection{Определение | Множество степени}
$P(A)$ --- множество подмножеств $A$, называемых собственными.

\subsubsection{Пример}
\begin{gather*}
    A = \{1, 2, 3\} \\
    P(A) = \{\varnothing, \{1\}, \{2\}, \{3\}, \{1, 2\}, \{1, 3\}, \{2, 3\}, A\}
\end{gather*}

\subsubsection{Пример}
\includegraphics[scale=0.2]{images/part01_01.png}

\subsubsection{Пример}
\includegraphics[scale=0.2]{images/part01_02.png}


\section{Операции над множествами}

\subsection{Определение | Объединение}
Объединение множеств $A$ и $B$
\[
    A \cup B = \{x~|~x \in A \text{ или } x \in B\}
\]
\includegraphics[scale=0.3]{images/part01_03.jpg}


\subsection{Определение | Пересечение}
Пересечение множеств $A$ и $B$
\[
    A \cap B = \{x~|~x \in A \text{ и } x \in B\}
\]
\includegraphics[scale=0.3]{images/part01_04.jpg}

\subsubsection{Пример (для обеих операций)}
\[
    A \subseteq B \Rightarrow A \cup B = B, A \cap B = A
\]
\includegraphics[scale=0.3]{images/part01_05.jpg}

\subsubsection{Пример (для обеих операций)}
Бесконечное множество концентрических квадратов (т.е. с общим центром) \\
\includegraphics[scale=0.2]{images/part01_06.jpg}


\subsection{Определение | Относительное дополнение (разность)}
Относительное дополнение множества $B$ до множества $A$
\[
    B \setminus A = \{x~|~x \in B \text{ и } x \notin A\}
\]
\includegraphics[scale=0.3]{images/part01_07.jpg}

\subsubsection{Пример}
\[
    A \subseteq B \Rightarrow A \setminus B = \varnothing
\] 


\subsection{Определение | Симметрическая разность}
Симметрическая разность множеств $A$ и $B$
\[
    A \oplus B = A + B = A \bigtriangleup B = A \setminus B \cup B \setminus A = (A \cup B) \setminus (A \cap B)
\]
\includegraphics[scale=0.3]{images/part01_08.jpg}


\subsection{Определение | Абсолютное дополнение (отрицание)}
Абсолютное дополнение множества $A$
\[
    \overline{A} = \{x~|~x \notin A\}
\]

\subsubsection{Пример}
\begin{gather*}
    B \setminus A = B \cap \overline{A} \\
    A \oplus B = (A \cap B) \cup (B \cap \overline{A}) = (A \cup B) \cap (\overline{A \cap B})
\end{gather*}

\subsubsection{Примечание}
$\mathbb{U}$ (или $\mathbb{E}$) --- универсальное множество (множество всех подмножеств). \\
\includegraphics[scale=0.3]{images/part01_09.jpg}

\subsection{Пример (для всех свойств)}
\paragraph{Дано}
\[
    A = [5 ; 9), B = [6, 10), C = [4; 8)
\]
\paragraph{Найти}
\[
    A \cup B;~~~A \cap C;~~~\overline{B};~~~C \setminus A;~~~B \oplus C;~~~\left( \overline{A} \setminus (B \cup \overline{C}) \right) \oplus \left( (\overline{B \setminus A}) \cup (\overline{C \setminus B}) \right)
\]
\paragraph{Решение}
\begin{enumerate}
    \item $A \cup B = [5; 10)$
    \item $A \cap C = [5; 8)$
    \item $\overline{B} = (-\infty; 6) \cup[10; +\infty)$
    \item $C \setminus A = [4; 5)$
    \item $B \oplus C = [4; 6) \cup [8; 10)$
    \item $\underbrace{\left( \underbrace{\underbrace{\overline{A}}_{\text{(a)}} \setminus \underbrace{(B \cup \underbrace{\overline{C}}_{\text{(b)}})}_{\text{(c)}}}_{\text{(d)}} \right) \oplus 
                       \left( \underbrace{\underbrace{(\overline{B \setminus A})}_{\text{(e, f)}} \cup \underbrace{(\overline{C \setminus B})}_{\text{(g, h)}}}_{\text{(i)}} \right)}_{\text{(j)}}$
    \begin{enumerate}
        \item $\overline{A} = (-\infty; 5) \cup [9; +\infty)$
        \item $\overline{C} = (-\infty; 4) \cup [8; +\infty)$
        \item $B \cup \overline{C} = (-\infty; 4) \cup [6; +\infty)$
        \item $\overline{A} \setminus (B \cup \overline{C}) = [4; 5)$
        \item $B \setminus A = [9; 10)$
        \item $\overline{B \setminus A} = (-\infty; -9) \cup [10; +\infty)$
        \item $C \setminus B = [4; 6)$
        \item $\overline{C \setminus B} = (-\infty; 4) \cup [6; +\infty)$
        \item $(\overline{B \setminus A}) \cup (\overline{C \setminus B}) = (-\infty; 4) \cup [6; +\infty)$
        \item $(-\infty; \infty) = \mathbb{R}$
    \end{enumerate}
\end{enumerate}



\section{Свойства операций (основные тождества)}

\subsection{Коммутативность (переместительный закон)}
\begin{gather*}
    A \cup B = B \cup A \\
    A \cap B = B \cap A \\
    A \oplus B = B \oplus A
\end{gather*}


\subsection{Ассоциативность (сочетательный закон)}
\begin{gather*}
    A \cup (B \cup C) = (A \cup B) \cup C = A \cup B \cup C \\
    A \cap (B \cap C) = (A \cap B) \cap C = A \cap B \cap C \\
    A \setminus (B \setminus C) = (A \setminus B) \setminus C
\end{gather*}


\subsection{Дистрибутивность (распределительный закон)}
\begin{gather*}
    A \cup (B \cap C) = (A \cup B) \cap (A \cup C) \\
    A \cap (B \cup C) = (A \cap B) \cup (A \cap C)
\end{gather*}

\subsubsection{Примечание}
Приоритет операций определяется только скобками.


\subsection{Закон идемпотентности}
\begin{gather*}
    A \cup A = A \\
    A \cap A = A
\end{gather*}


\subsection{Закон де Моргана}
\begin{gather*}
    \overline{A \cup B} = \overline{A} \cap \overline{B} \\
    \overline{A \cap B} = \overline{A} \cup \overline{B}
\end{gather*}


\subsection{Закон поглощения}
\begin{gather*}
    A \cup (A \cap B) = A \\
    A \cap (A \cup B) = A
\end{gather*}


\subsection{Расщепления}
\begin{gather*}
    A = (A \cup B) \cap (A \cup \overline{B})
\end{gather*}


\subsection{Двойное отрицание}
\begin{gather*}
    \overline{\overline{A}} = A \\
    X \setminus Y = X \cap \overline{Y} \\
    X \oplus Y = (X \cap \overline{Y}) \cup (\overline{X} \cap Y) = (X \cup Y) \cap (\overline{X} \cup \overline{Y})
\end{gather*}


\subsection{Разности}
\begin{gather*}
    A \setminus (B \cup C) = (A \setminus B) \cap (A \setminus C) \\
    A \setminus (B \cap C) = (A \setminus B) \cup (A \setminus C) \\ 
    (A \setminus B) \cap C = (A \cap C) \setminus (B \cap C) \\
    (A \setminus B) \setminus C = (A \setminus C) \setminus B
\end{gather*}



\section{Характеристическая функция множества}

\subsection{Определение}
\begin{gather*}
    A \subseteq \mathbb{U} \\
    A \Rightarrow \psi_A: \mathbb{U} \Rightarrow \{0, 1\} \\
    \psi_A(x) = 
    \begin{cases}
        1, \text{ если } x \in A \\
        0, \text{ если } x \notin A
    \end{cases}
\end{gather*}


\subsection{Функции основных операций через характеристические функции}
$\psi_A(x), \psi_B(x)$
\begin{enumerate}
    \item $\psi_{A \cap B} = \psi_A \cdot \psi_B$
    \item $\psi_{A \cup B} = \psi_A + \psi_B - \psi_A \cdot \psi_B$
    \item $\psi_{\overline{A}} = 1 - \psi_A$
    \item $\psi_{A \setminus B} = \psi_{A \cap \overline{B}} = \psi_A \cdot (1 - \psi_B) = \psi_A - \psi_A \psi_B$
    \item $\psi^2_A = \psi_A$
    \item $\psi_{A \oplus B} = \psi_A + \psi_B - 2 \psi_A \psi_B = \psi^2_A + \psi^2_B - 2 \psi_A \psi_B = (\psi_A - \psi_B)^2 = |\psi_A - \psi_B|$
\end{enumerate}

\subsubsection{Пример}
Доказать, что
\[
    A \setminus (B \setminus C) = (A \setminus B) \cup (A \cap C)
\]

\paragraph{Доказательство}
\begin{gather*}
    \psi_{A \setminus (B \setminus C)} = 
    \psi_A \cdot (1 - \psi_{B \setminus C}) = 
    \psi_A (1 - \psi_B (1 - \psi_C)) =
    \psi_A - \psi_A \psi_B + \psi_A \psi_B \psi_C \\
    \psi_{(A \setminus B) \cup (A \cap C)} = 
    \psi_A (1 - \psi_B) + \psi_A \psi_C - \psi_A (1 - \psi_B) \cdot \psi_A \psi_C = \\
    = \psi_A - \psi_A \psi_B + \psi_A \psi_C - \psi_A \psi_C + \psi_A \psi_B \psi_C = \psi_A - \psi_A \psi_B + \psi_A \psi_B \psi_C
\end{gather*}


\subsubsection{Пример}
Доказать, что
\[
    A \cap (B \oplus C) = (A \cap B) \oplus (A \cap C)
\]

\paragraph{Доказательство}
\begin{gather*}
    \psi_{A \cap (B \oplus C)} = \psi_A \psi_{B \oplus C} = \psi_A |\psi_B - \psi_C| \\
    \psi_{(A \cap B) \oplus (A \cap C)} = |\psi_{A \cap B} - \psi_{A \cap C}| = \psi_A |\psi_B - \psi_C|
\end{gather*}


\section{Таблица характеристических операций}


\subsection{Определение}
\begin{tabular}{L | L | L | L | L | L | L}
    \psi_A & \psi_B & \psi_{A \cap B} & \psi_{A \cup B} & \psi_{\overline{A}} & \psi_{A \setminus B} & \psi_{A \oplus B} \\
    \hline
    1 & 1 & 1 & 1 & 0 & 0 & 0 \\
    1 & 0 & 0 & 1 & 0 & 1 & 1 \\
    0 & 1 & 0 & 1 & 1 & 0 & 1 \\
    0 & 0 & 0 & 0 & 1 & 0 & 0
\end{tabular}


\subsection{Практика}

\subsubsection{Пример}
Доказать, что
\[
    A \setminus (B \setminus C) = (A \setminus B) \cup (A \cap C)
\]

\paragraph{Доказательство}
~\\
\begin{tabular}{L | L | L | L | L || L | L | L}
    \psi_{A} & \psi_{B} & \psi_{С} & \psi_{B \setminus C} & \psi_{A \setminus (B \setminus C)} & \psi_{A \setminus B} & \psi_{A \cap C} & \psi_{(A \setminus B) \cup (A \cap C)} \\
    \hline
    1 & 1 & 1 & 0 & \textcolor{green}{1} & 0 & 1 & \textcolor{green}{1} \\
    1 & 1 & 0 & 1 & \textcolor{green}{0} & 0 & 0 & \textcolor{green}{0} \\
    1 & 0 & 1 & 0 & \textcolor{green}{1} & 1 & 1 & \textcolor{green}{1} \\
    1 & 0 & 0 & 0 & \textcolor{green}{1} & 0 & 0 & \textcolor{green}{1} \\
    0 & 1 & 1 & 0 & \textcolor{green}{0} & 0 & 0 & \textcolor{green}{0} \\
    0 & 1 & 0 & 1 & \textcolor{green}{0} & 0 & 0 & \textcolor{green}{0} \\
    0 & 0 & 1 & 0 & \textcolor{green}{0} & 0 & 0 & \textcolor{green}{0} \\
    0 & 0 & 0 & 0 & \textcolor{green}{0} & 0 & 0 & \textcolor{green}{0}
\end{tabular}


\subsubsection{Пример}
Доказать, что
\[
    A \cap (B \oplus C) = (A \cap B) \oplus (A \cap C)
\]

\paragraph{Доказательство}
~\\
\begin{tabular}{L | L | L | L | L | L | L | L}
    \psi_{A} & \psi_{B} & \psi_{С} & \psi_{B \oplus C} & \psi_{A \cap (B \oplus C)} & \psi_{A \cap B} & \psi_{A \cap C} & \psi_{A \cap C} \\
    \hline
    1 & 1 & 1 & 0 & \textcolor{green}{0} & 1 & 1 & \textcolor{green}{0} \\
    1 & 1 & 0 & 1 & \textcolor{green}{1} & 1 & 0 & \textcolor{green}{1} \\
    1 & 0 & 1 & 1 & \textcolor{green}{1} & 0 & 1 & \textcolor{green}{1} \\
    1 & 0 & 0 & 0 & \textcolor{green}{0} & 0 & 0 & \textcolor{green}{0} \\
    0 & 1 & 1 & 0 & \textcolor{green}{0} & 0 & 0 & \textcolor{green}{0} \\
    0 & 1 & 0 & 1 & \textcolor{green}{0} & 0 & 0 & \textcolor{green}{0} \\
    0 & 0 & 1 & 1 & \textcolor{green}{0} & 0 & 0 & \textcolor{green}{0} \\
    0 & 0 & 0 & 0 & \textcolor{green}{0} & 0 & 0 & \textcolor{green}{0}
\end{tabular}



\section{Методы утверждений}
\begin{enumerate}
    \item Теорема --- важный результат
    \item Утверждение --- теорема с менее важным результатом
    \item Лемма --- вспомогательное утверждение для теоремы
\end{enumerate}


\subsection{Утверждение}
3 эквивалентных условия (выполняются все при выполнении одного)
\begin{enumerate}
    \item $A \subseteq B$
    \item $A \cap B = A$
    \item $A \cup B = B$
\end{enumerate}

\subsubsection{Доказательство 1 $\Rightarrow$ 2}
\paragraph{Дано}
\[
    A \subseteq B
\]

\paragraph{Доказать}
\[
    A \cap B = A
\]

\paragraph{Доказательство}
\begin{gather*}
    x = y \Leftrightarrow x \subseteq y, y \subseteq x \\
    x \subseteq y: \forall x \in X \Rightarrow x \in y
\end{gather*}
\begin{enumerate}
    \item $A \cap B \subseteq A$
    \[
        \text{Пусть } x \in A \cap B \Rightarrow x \in A.
    \]
    \item $A \subseteq A \cap B$
    \[
        \text{Пусть } x \in A \underset{\text{усл.}}{\Rightarrow} x \in B \Rightarrow x \in A \cap B.
    \]
\end{enumerate}


\subsubsection{Доказательство 2 $\Rightarrow$ 3}
\paragraph{Дано}
\[
    A \cap B = A
\]

\paragraph{Доказать}
\[
    A \cup B = B
\]

\paragraph{Доказательство}
\[
    A \cup B = (A \cap B) \cup B = B \text{ --- поглощение}
\]


\subsubsection{Доказательство 3 $\Rightarrow$ 1}
\paragraph{Дано}
\[
    A \cup B = B
\]

\paragraph{Доказать}
\[
    A \subseteq B
\]

\paragraph{Доказательство (от противного)}
\begin{gather*}
    \text{Предположим, что } A \nsubseteq B \Rightarrow \exists x \in A \text{ и } x \notin B. \\ 
    \text{Так как } x \in A \Rightarrow x \in A \cup B \underset{\text{усл.}}{\Rightarrow} x \in B \text{ --- противоречие} \Rightarrow A \subseteq B
\end{gather*}




\newpage
\paragraph{2024.10.01 (3, 4)}
\part{Методы доказательств (опровержений) тождеств}

\section{Метод рассуждений}

\subsection{Пример | И. А. Лавров, Л. Л. Максимова: номер 13(г)}
Доказать, что 
\[
    A \subseteq B \cup C \Leftrightarrow A \cap \overline{B} \subseteq C
\]

\subsubsection{Доказательство 1 $\Rightarrow$ 2}
\paragraph{Дано}
\[
    A \subseteq B \cup C
\]

\paragraph{Доказать}
\[
    A \cap \overline{B} \subseteq C
\]

\paragraph{Доказательство}
\begin{gather*}
    \text{Пусть } x \in A \cap \overline{B} \Rightarrow x \in A \text{ и } x \in \overline B. \\
    x \in A \underset{\text{усл.}}{\Rightarrow} x \in B \cup C \Rightarrow x \in B \text{ или } x \in C, \text{ но } x \notin B \Rightarrow x \in C.
\end{gather*}


\subsubsection{Доказательство 1 $\Leftarrow$ 2}
\paragraph{Дано}
\[
    A \cap \overline{B} \subseteq C
\]

\paragraph{Доказать}
\[
    A \subseteq B \cup C
\]

\paragraph{Доказательство}
\begin{gather*}
    \text{Пусть } x \in A. \\
    \text{Если } x \in B \Rightarrow x \in B \cup C. \\
    \text{Если } x \notin B \Rightarrow x \in A \cap \overline B \underset{\text{усл.}}{\Rightarrow} x \in C \Rightarrow x \in B \cup C.
\end{gather*}


\subsection{Пример}
\begin{gather*}
    X = Y \Rightarrow X \cup Z = Y \cup Z \\
    X \cup Z = Y \cup Z \nRightarrow X = Y
\end{gather*}


\subsection{Пример}
Доказать, что
\[
    A \oplus B = B \setminus C \Leftrightarrow A = B \cap C
\]

\subsubsection{Доказательство 1 $\Leftarrow$ 2}
\paragraph{Дано}
\[
    A = B \cap C
\]

\paragraph{Доказать}
\[
    A \oplus B = B \setminus C
\]

\paragraph{Доказательство}
\begin{gather*}
    (B \cap C) \oplus B = (B \cap C \cap \overline{B}) \cup \left( (\overline{B \cap C}) \cap B \right) = 
    (\overline{B} \cup \overline{C}) \cap B = (\overline{B} \cap B) \cup (B \cap C) = B \cap C
\end{gather*}


\subsubsection{Доказательство 1 $\Rightarrow$ 2}
\paragraph{Дано}
\[
    A \oplus B = B \setminus C
\]

\paragraph{Доказать}
\[
    A = B \cap C
\]

\paragraph{Доказательство 1}
\begin{enumerate}
    \item Докажем, что $A \subseteq B \cap C$
    \begin{gather*}
        \text{Пусть } x \in A. \\
        \text{Если } x \in B \Rightarrow x \notin A \oplus B \underset{\text{усл.}}{\Rightarrow} 
        x \notin B \setminus C \Rightarrow x \in \overline{B \setminus C} = \overline{B \cap \overline{C}} = \overline{B} \cup C \Rightarrow \\
        x \in \overline{B} \text{ или } x \in C, \text{ но } x \in B \text{ --- противоречие } \Rightarrow x \in B.
    \end{gather*}
    \item Докажем, что $B \cap C \subseteq A$
    \begin{gather*}
    \text{Пусть } x \in B \cap C \Rightarrow x \in B \text{ и } x \in C \Rightarrow 
    x \notin B \setminus C \underset{\text{усл.}}{\Rightarrow} x \notin A \oplus B \Rightarrow x \in A.
    \end{gather*}
\end{enumerate}

\paragraph{Доказательство 2}
\begin{gather*}
    (A \oplus B) \oplus B = (B \setminus C) \oplus B \\
    A = (B \setminus C) \oplus B \\
    (B \setminus C) \oplus B = 
    (B \cap \overline{C} \cap \overline{B}) \cup \left( (\overline{B \cap \overline{C}}) \cap B \right) = 
    (\overline{B} \cup C) \cap B = 
    (B \cap \overline{B}) \cup (B \cap C) = B \cap C.
\end{gather*}


\subsection{Пример | И. А. Лавров, Л. Л. Максимова: номер 13(е)}
Доказать, что
\[
    (A \cap B) \cup C = A \cap (B \cup C) \Leftrightarrow C \subseteq A
\]

\subsubsection{Доказательство 1 $\Leftarrow$ 2}
\paragraph{Дано}
\[
C \subseteq A
\]

\paragraph{Доказать}
\[
    (A \cap B) \cup C = A \cap (B \cup C)
\]

\paragraph{Доказательство}
\begin{gather*}
    C \subseteq A \Rightarrow C \cup A = A, C \cap A = C. \\
    (A \cap B) \cup C = (A \cup C) \cap (B \cup C) = A \cap (B \cup C)
\end{gather*}


\subsubsection{Доказательство 1 $\Rightarrow$ 2}
\paragraph{Дано}
\[
    (A \cap B) \cup C = A \cap (B \cup C)
\]

\paragraph{Доказать}
\[
    C \subseteq A
\]

\paragraph{Доказательство}
\begin{gather*}
    \text{Пусть } x \in C. \\
    \text{Если } x \in C \Rightarrow (A \cap B) \cup C \Rightarrow x \in A \cap (B \cup C) \Rightarrow x \in A
\end{gather*}


\subsection{Пример}
Доказать, что
\[
    A \subseteq B \Leftrightarrow A \setminus B = \varnothing
\]

\subsubsection{Доказательство 1 $\Rightarrow$ 2}
\paragraph{Дано}
\[
    A \subseteq B
\]

\paragraph{Доказать}
\[
    A \setminus B = \varnothing
\]

\paragraph{Доказательство (от противного)}
\begin{gather*}
    \text{Пусть } A \setminus B \neq \varnothing \Rightarrow x \in A \setminus B \Rightarrow \\
    \Rightarrow x \in A \text{ и } x \notin B, \text{ но } A \subseteq B \text{ --- противоречие } \Rightarrow A \setminus B = \varnothing
\end{gather*}


\subsubsection{Доказательство 1 $\Leftarrow$ 2}
\paragraph{Дано}
\[
    A \setminus B = \varnothing
\]

\paragraph{Доказать}
\[
    A \subseteq B
\]

\paragraph{Доказательство (от противного)}
\begin{gather*}
    \text{Пусть } A \nsubseteq B \Rightarrow x \in A \text{ и } x \notin B \Rightarrow \\
    \Rightarrow x \in A \setminus B, \text{ но } A \setminus B = \varnothing \text{ --- противоречие} \Rightarrow A \subseteq B
\end{gather*}


\subsection{Пример}
Доказать, что
\[
    A = B \Leftrightarrow A \oplus B = \varnothing
\]

\subsubsection{Доказательство 1 $\Leftarrow$ 2}
\paragraph{Дано}
\[
    A \oplus B = \varnothing
\]

\paragraph{Доказать}
\[
    A = B
\]

\paragraph{Доказательство}
\begin{gather*}
    A \oplus B = \varnothing \Leftrightarrow (A \setminus B) \cup (B \setminus A) = \varnothing \Leftrightarrow
    \begin{cases}
        A \setminus B = \varnothing \\
        B \setminus A = \varnothing
    \end{cases} \Leftrightarrow
    \begin{cases}
        A \subseteq B \\
        B \subseteq A
    \end{cases} \Leftrightarrow
    A = B
\end{gather*}



\section{Характеристическая функция}

\subsection{Утверждение}
Пусть $F_1$, $F_2$, $G_1$, $G_2$ --- формулы, тогда
\begin{enumerate}
    \item $F_1 = F_2 \Leftrightarrow G_1 = G_2$ выполняется тогда и только тогда, когда $\psi_{F_1 \oplus F_2} = \psi_{G_1 \oplus G_2}$.
    \item $F_1 = F_2 \Leftrightarrow G_1 \subseteq G_2$ выполняется тогда и только тогда, когда $\psi_{F_1 \oplus F_2} = \psi_{G_1 \setminus G_2}$.
    \item $F_1 \subseteq F_2 \Leftrightarrow G_1 = G_2$ выполняется тогда и только тогда, когда $\psi_{F_1 \setminus F_2} = \psi_{G_1 \oplus G_2}$.
    \item $F_1 \subseteq F_2 \Leftrightarrow G_1 \subseteq G_2$ выполняется тогда и только тогда, когда $\psi_{F_1 \setminus F_2} = \psi_{G_1 \setminus G_2}$.
\end{enumerate}


\subsection{Практика}

\subsubsection{Пример | И. А. Лавров, Л. Л. Максимова: номер 13(г)}
Доказать, что
\[
    A \subseteq B \cup C \Leftrightarrow A \cap \overline{B} \subseteq C
\]

\paragraph{Доказательство}
\begin{gather*}
    A \setminus (B \cup C) = \varnothing \Leftrightarrow (A \cap \overline{B}) \setminus C = \varnothing \\
    \psi_{A \setminus (B \cup C)} = 0 \Leftrightarrow \psi_{(A \cap \overline{B}) \setminus C} = 0 \\
    \psi_{A \setminus (B \cup C)} = \psi_{(A \cap \overline{B}) \setminus C} \\
    \begin{tabular}{c|c|c|c|c|c|c|c}
        $\psi_A$ & $\psi_B$ & $\psi_C$ & $\psi_{B \cup C}$ & $\psi_{A \setminus (B \cup C)}$ & $\psi_{\overline{B}}$ & $\psi_{A \cap \overline{B}}$ & $\psi_{(A \cap \overline{B}) \setminus C}$ \\
        \hline
        0 & 0 & 0 & 0 & \textcolor{green}{0} & 1 & 0 & \textcolor{green}{0} \\
        0 & 0 & 1 & 1 & \textcolor{green}{0} & 1 & 0 & \textcolor{green}{0} \\
        0 & 1 & 0 & 1 & \textcolor{green}{0} & 0 & 0 & \textcolor{green}{0} \\
        0 & 1 & 1 & 1 & \textcolor{green}{0} & 0 & 0 & \textcolor{green}{0} \\
        1 & 0 & 0 & 0 & \textcolor{green}{1} & 1 & 1 & \textcolor{green}{1} \\
        1 & 0 & 1 & 1 & \textcolor{green}{0} & 1 & 1 & \textcolor{green}{0} \\
        1 & 1 & 0 & 1 & \textcolor{green}{0} & 0 & 0 & \textcolor{green}{0} \\
        1 & 1 & 1 & 1 & \textcolor{green}{0} & 0 & 0 & \textcolor{green}{0} \\  
    \end{tabular}
\end{gather*}


\subsubsection{Пример}
Доказать, что
\[
    A \cup B \subseteq C \Leftrightarrow A \subseteq C \text{ и } B \subseteq C
\]

\paragraph{Доказательство}
\begin{gather*}
    (A \cup B) \setminus C = \varnothing \Leftrightarrow A \setminus C = \varnothing \text{ и } B \setminus C = \varnothing \\
    \psi_{(A \cup B) \setminus C} = 0 \Leftrightarrow \psi_{A \setminus C} \text{ и } \psi_{B \setminus C} = 0 \\
    \begin{tabular}{c|c|c|c|c||c|c}
        $\psi_A$ & $\psi_B$ & $\psi_C$ & $\psi_{A \cup B}$ & $\psi_{(A \cup B) \setminus C}$ & $\psi_{A \setminus C}$ & $\psi_{B \setminus C}$ \\
        \hline
        0 & 0 & 0 & 0 & \textcolor{green}{0}              & \textcolor{green}{0} & \textcolor{green}{0} \\
        0 & 0 & 1 & 0 & \textcolor{green}{0}              & \textcolor{green}{0} & \textcolor{green}{0} \\
        0 & 1 & 0 & 1 & \textcolor{red}{1} $\nLeftarrow$  & \textcolor{red}{0}   & \textcolor{red}{1}   \\
        0 & 1 & 1 & 1 & \textcolor{green}{0}              & \textcolor{green}{0} & \textcolor{green}{0} \\
        1 & 0 & 0 & 1 & \textcolor{green}{1}              & \textcolor{green}{1} & \textcolor{green}{0} \\
        1 & 0 & 1 & 1 & \textcolor{red}{0} $\nRightarrow$ & \textcolor{red}{0}   & \textcolor{red}{0}   \\
        1 & 1 & 0 & 1 & \textcolor{green}{1}              & \textcolor{green}{1} & \textcolor{green}{1} \\
        1 & 1 & 1 & 1 & \textcolor{green}{0}              & \textcolor{green}{0} & \textcolor{green}{0} \\  
    \end{tabular}
\end{gather*}


\subsubsection{Пример}
Доказать, что
\[
    A \cup B \subseteq C \Leftrightarrow A \subseteq C \text{ и } A \subseteq B
\]

\paragraph{Доказательство}
\begin{gather*}
    (A \cup B) \setminus C = \varnothing \Leftrightarrow A \setminus C = \varnothing \text{ и } A \setminus B = \varnothing \\
    \psi_{(A \cup B) \setminus C} = 0 \Leftrightarrow \psi_{A \setminus C} \text{ и } \psi_{A \setminus B} = 0 \\
    \begin{tabular}{c|c|c|c|c||c|c}
        $\psi_A$ & $\psi_B$ & $\psi_C$ & $\psi_{A \cup B}$ & $\psi_{(A \cup B) \setminus C}$ & $\psi_{A \setminus C}$ & $\psi_{A \setminus B}$ \\
        \hline
        0 & 0 & 0 & 0 & \textcolor{green}{0}              & \textcolor{green}{0} & \textcolor{green}{0} \\
        0 & 0 & 1 & 0 & \textcolor{green}{0}              & \textcolor{green}{0} & \textcolor{green}{0} \\
        0 & 1 & 0 & 1 & \textcolor{red}{1} $\nLeftarrow$  & \textcolor{red}{0}   & \textcolor{red}{1}   \\
        0 & 1 & 1 & 1 & \textcolor{green}{0}              & \textcolor{green}{0} & \textcolor{green}{0} \\
        1 & 0 & 0 & 1 & \textcolor{green}{1}              & \textcolor{green}{1} & \textcolor{green}{0} \\
        1 & 0 & 1 & 1 & \textcolor{red}{0} $\nRightarrow$ & \textcolor{red}{0}   & \textcolor{red}{0}   \\
        1 & 1 & 0 & 1 & \textcolor{green}{1}              & \textcolor{green}{1} & \textcolor{green}{1} \\
        1 & 1 & 1 & 1 & \textcolor{green}{0}              & \textcolor{green}{0} & \textcolor{green}{0} \\  
    \end{tabular}
\end{gather*}


\subsubsection{Пример}
Доказать, что
\[
    A \oplus B = B \setminus C \Leftrightarrow A = B \cap C
\]

\paragraph{Доказательство}
\begin{gather*}
    \psi_{(A \oplus B) \oplus (B \setminus C)} = \psi_{A \oplus (B \cap C)} \\
    \begin{tabular}{c|c|c|c|c|c|c|c}
        $\psi_A$ & $\psi_B$ & $\psi_C$ & $\psi_{A \oplus B}$ & $\psi_{B \setminus C}$ & $\psi_{(A \oplus B) \oplus (B \setminus C)}$ & $\psi_{B \cap C}$ & $\psi_{A \oplus (B \cap C)}$ \\
        \hline
        0 & 0 & 0 & 0 & 0 & \textcolor{green}{0} & 0 & \textcolor{green}{0} \\
        0 & 0 & 1 & 0 & 0 & \textcolor{green}{0} & 0 & \textcolor{green}{0} \\
        0 & 1 & 0 & 1 & 1 & \textcolor{green}{0} & 0 & \textcolor{green}{0} \\
        0 & 1 & 1 & 1 & 0 & \textcolor{green}{1} & 1 & \textcolor{green}{1} \\
        1 & 0 & 0 & 1 & 0 & \textcolor{green}{1} & 0 & \textcolor{green}{1} \\
        1 & 0 & 1 & 1 & 0 & \textcolor{green}{1} & 0 & \textcolor{green}{1} \\
        1 & 1 & 0 & 0 & 1 & \textcolor{green}{1} & 0 & \textcolor{green}{1} \\
        1 & 1 & 1 & 0 & 0 & \textcolor{green}{0} & 1 & \textcolor{green}{0} \\  
    \end{tabular}
\end{gather*}



\section{Решение уравнений и системы уравнений}
Пусть $A$, $B$, $C \subseteq \mathbb{U}$, $X$ --- ? ~~~~~ $\oplus$ --- $+$, $\varnothing$ --- $0$

\subsection{Практика}

\subsubsection{Пример}
\begin{gather*}
    A + X = 0 \\
    X = A
\end{gather*}

\subsubsection{Пример}
\begin{gather*}
    A + X = B~~~|+B \\
    A + B + X = 0 \\
    X = A + B
\end{gather*}

\subsubsection{Пример}
\begin{gather*}
    A \cap X = 0 \\
    X \setminus \overline{A} = 0 \\
    X \subseteq \overline{A}
\end{gather*}
\includegraphics[scale=0.4]{images/part02_01.png}

\subsubsection{Пример}
\[
    A \cap X = B
\]

\paragraph{Решение}
\begin{gather*}
    (A \cap X) + B = \varnothing \\
    (A \cap X \cap \overline{B}) \cup \left( (\overline{A \cap X}) \cap B \right) = \varnothing \\
    (A \cap X \cap \overline{B}) \cup \left( (\overline{A} \cup \overline{X}) \cap B \right) = \varnothing \\
    (X \cap A \cap \overline{B}) \cup (\overline{A} \cap B) \cup (\overline{X} \cap B) = \varnothing \\
    \begin{cases}
        X \cap A \cap \overline{B} = \varnothing ~~~ (1) \\
        \overline{A} \cap B = \varnothing                \\
        \overline{X} \cap B = \varnothing ~~~~~~~~~  (3)
    \end{cases} \Leftrightarrow
    \begin{cases}
        B \subseteq A \\
        B \subseteq X
    \end{cases} \\
    X \cap (A \cap \overline{B}) = X \setminus (\overline{A \cap \overline{B}}) = X \setminus (\overline{A} \cup B) = \varnothing \\
    (1) ~~~ (A \cup B) \cap \overline{C} \cap A \cap \overline{B} = 
    (\overline{A} \cap \overline{C} \cap A \cap B) \cup (B \cap \overline{C} \cap A \cap \overline{B}) = \varnothing \\
    (3) ~~~ \left( (\overline{\overline{A} \cup B}) \cap \overline{C} \right) \cap B =
    \left( (A \cap B) \cup C \right) \cap B =
    (A \cap \overline{B} \cap B) \cup (C \cap B) = B \cap C = \varnothing \\
    B \subseteq A, C \subseteq \overline{A}
\end{gather*}
\includegraphics[scale=0.3]{images/part02_02.png}


\paragraph{Ответ}
\[
    X = (\overline{A} \cup B) \setminus C, B \subseteq A, C \subseteq \overline{A}
\]


\subsubsection{Пример}
\begin{gather*}
\begin{cases}
    A \cup X = B \\
    A \cup X = C
\end{cases}
\end{gather*}

\paragraph{Решение}
\begin{gather*}
    (A \cup X) \oplus C = \varnothing; \\
    \left( (A \cup X) \cap \overline{C} \right) \cup \left( (\overline{A \cup X}) \cap C \right) = \varnothing \\
    (A \cap \overline{C}) \cup (X \cap \overline{C}) \cup (\overline{A} \cap \overline{X} \cap C) = \varnothing \\
    \begin{cases}{}
        X \cap A \cap \overline{B} = \varnothing \\
        \overline{A} \cap B = \varnothing \\
        \overline{X} \cap B = \varnothing \\
        A \cap \overline{C} = \varnothing \\
        X \cap \overline{C} = \varnothing \\
        \overline{A} \cap \overline{X} \cap \overline{C} = \varnothing
    \end{cases} \Leftrightarrow
    \begin{cases}{}
        X \subseteq \overline{A} \cup B ~~~ (1) \\
        B \subseteq A ~~~~~~~~~ (2) \\
        B \subseteq X ~~~~~~~~~ (3) \\
        A \subseteq C ~~~~~~~~~ (4) \\
        X \subseteq C ~~~~~~~~~ (5) \\
        \overline{A} \cap C \subseteq X ~~~ (6) \\
    \end{cases}
\end{gather*}
\includegraphics[scale=0.6]{images/part02_03.png}
\begin{gather*}
    (3), (6):
    \begin{cases}
        B \subseteq X \\
        \overline{A} \cap C \subseteq X
    \end{cases} \Leftrightarrow
    B \cup (\overline{A} \cap C) = X \\
    (1), (5):
    \begin{cases}
        X \subseteq \overline{A} \cup B \\
        X \subseteq C
    \end{cases} \Leftrightarrow
    X \subseteq (\overline{A} \cup B) \cap C \\
    X \subseteq (\overline{A} \cap C) \cup (B \cap C) \Leftrightarrow X \subseteq (A \cap C) \cup B
\end{gather*}



\newpage
\paragraph{2024.10.15 (3, 4)}
\part{Бинарные отношения}

\section{Введение}

\subsection{Определение | Упорядоченная пара}
\begin{gather*}
    \langle x, y \rangle \\
    \langle x, y \rangle = \langle a, b \rangle \Leftrightarrow x = a, b = y \\
\end{gather*}


\subsection{Определение | Бинарные отношения $\rho$ на множестве $X$}
\begin{gather*}
    \rho = \{\langle x, y \rangle | ~x, y \in X \rho(x, y)\} \\
    \langle x, y \rangle \in \rho, x \rho y \\
    \langle x, y \rangle \notin \rho, x \cancel{\rho} y
\end{gather*}

\subsubsection{Пример}
\begin{gather*}
    \rho \text{ --- } =, \mathbb{R}:~ \langle 2, 2 \rangle \in =, 2 = 2 \\
    \rho \text{ --- } \le, \mathbb{R}:~ \langle 2, 5 \rangle \in \le, 2 \le 5
\end{gather*}



\section{Характеристики}

\subsection{Область определения $\rho$ на $X$}
\[
D_{\rho} = \{x \in X ~|~ x \rho y\}
\]


\subsection{Область значения $\rho$ на $X$}
\[
R_{\rho} = \{y \in X ~|~ x \rho y\}
\]

\subsection{Пример}
\begin{gather*}
    \rho = \{\langle x, y \rangle |~\mathbb{R}, y = \sqrt{x - 1}\} \\
    D_{\rho} = [1; +\infty) \\
    R_{\rho} = (0; +\infty) \\
    \\
    X = \{1, 2, 3\} \\
    \rho = \{\langle 1, 2 \rangle, \langle 1, 3 \rangle, \langle 2, 3 \rangle\} \\
    D_{\rho} = \{1, 2\}; R_{\rho} = \{2, 3\}
\end{gather*}



\section{Прямое (декартовое) произведение}

\subsection{Определение}
\[
X, Y~~~X \times Y = \{\langle x, y \rangle |~x \in X, y \in Y\}
\]

\subsubsection{Пример}
\begin{gather*}
    X = \{1, 2, 3\};~Y = \{4, 5\} \\
    X \times Y = \{\langle 1, 4 \rangle, \langle 1, 5 \rangle, \langle 2, 4 \rangle, \langle 2, 5 \rangle, \langle 3, 4 \rangle, \langle 3, 5 \rangle\} \\
    Y \times X = \{\langle 4, 1 \rangle, \langle 4, 2 \rangle, \langle 4, 3 \rangle, \langle 5, 1 \rangle, \langle 5, 2 \rangle, \langle 5, 3 \rangle\} \\
\end{gather*}
\includegraphics[scale=0.3]{images/part03_01.jpg}


\subsection{Свойства}
\begin{gather*}
    X \times Y \neq Y \times X \\
    X \times X = X^2 \\
    D_{X \times Y} = X \\
    R_{X \times Y} = Y
\end{gather*}

\subsubsection{Пример}
\[
R^3 = R \times R \times R
\]
\includegraphics[scale=0.3]{images/part03_02.jpg}


\subsubsection{Пример}
\[
    (A \times B) \setminus (C \times D) = (A \times (B \setminus D)) \cup ((A \setminus C) \times B) \\
\]
\begin{enumerate}
    \item Докажем, что $(A \times B) \ (C \times D) \subseteq (A \times (B \setminus D)) \cup ((A \setminus C) \times B)$
    \begin{gather*}
        \text{Пусть } \langle x, y \rangle \in (A \times B) \setminus (C \times D) \Rightarrow \langle x, y \rangle \in A \times B \text{ и } \langle x, y \rangle \notin C \times D \Rightarrow \\
        \\
        \Rightarrow \begin{cases}
            x \in A \\
            y \in B
        \end{cases} \text{ и } (x \notin C \text{ или } y \notin D) \\
        \\
        \text{Если } \begin{cases}
            x \notin C, x \in A \\
            y \in B
        \end{cases}
        \Rightarrow \begin{cases}
            x \in A \setminus C \\
            y \in B
        \end{cases} \Rightarrow \\
        \Rightarrow \langle x, y \rangle \in (A \setminus C) \times B \Rightarrow \langle x, y \rangle \in (A \times (B \setminus D)) \cup ((A \setminus C) \times B). \\
        \\
        \text{Если } \begin{cases}
            x \in A \\
            y \notin D, y \in B
        \end{cases}
        \Rightarrow \begin{cases}
            x \in A \\
            y \in B \setminus D
        \end{cases} \Rightarrow \\
        \Rightarrow \langle x, y \rangle \in A \times (B \setminus D) \Rightarrow \langle x, y \rangle \in (A \times (B \setminus D)) \cup ((A \setminus C) \times B). \\
    \end{gather*}
    \item Докажем, что $(A \times (B \setminus D)) \cup ((A \setminus C) \times B) \subseteq (A \times B) \setminus (C \times D)$ \\
    \begin{gather*}
        \text{Пусть } \langle x, y \rangle \in (A \times (B \setminus D)) \cup ((A \setminus C) \times B) \Rightarrow \\ \langle x, y \rangle \in A \times (B \setminus D) \text{ или } \langle x, y \rangle \in (A \setminus C) \times B. \\
        \text{Если } \langle x, y \rangle \in A \times (B \setminus D) \Rightarrow
        \begin{cases}
            x \in A \\
            y \in B \setminus D
        \end{cases}
        \Rightarrow \begin{cases}
            x \in A \\
            y \in B, y \notin D
        \end{cases} \Rightarrow \\
        \Rightarrow \langle x, y \rangle \in A \times B, \langle x, y \rangle \notin C \times D \text { (так как } y \notin D) \Rightarrow \langle x, y \rangle \in (A \times B) \setminus (C \times D). \\
        \text{Если } \langle x, y \rangle \in (A \setminus C) \times B \Rightarrow
        \begin{cases}
            x \in A \setminus C \\
            y \in B
        \end{cases}
        \Rightarrow \begin{cases}
            x \in A, x \notin C \\
            y \in B
        \end{cases} \Rightarrow \\
        \Rightarrow \langle x, y \rangle \in A \times B, \langle x, y \rangle \notin C \times D \text{ (так как } x \notin C) \Rightarrow \langle x, y \rangle \in (A \times B) \setminus (C \times D).
    \end{gather*}
\end{enumerate}


\subsubsection{Пример}
\[
    L = (A \times B) \cap (C \times D) = (A \cap C) \times (B \cap D) = R
\]
\begin{enumerate}
    \item Докажем, что $L \subseteq R$
    \begin{gather*}
        \text{Пусть } \langle x, y \rangle \in L \Rightarrow
        \begin{cases}
            x \in A, x \in C \\
            y \in B, y \in D
        \end{cases} \Rightarrow 
        \begin{cases}
            x \in A \cap C \\
            y \in B \cap D
        \end{cases}
        \Rightarrow \langle x, y \rangle \in R. \\
    \end{gather*}
    \item Докажем, что $R \subseteq L$
    \begin{gather*}
        \text{Пусть } \langle x, y \rangle \in R \Rightarrow
        \begin{cases}
            x \in A \cap C \\
            y \in B \cap D
        \end{cases} \Rightarrow 
        \begin{cases}
            x \in A, x \in C \\
            y \in B, y \in D
        \end{cases}
        \Rightarrow \langle x, y \rangle \in L.
    \end{gather*}
\end{enumerate}


\section{Операции над бинарными отношениями}

\subsection{???}
\[
    \cup, \cap, \dots
\]


\subsection{Абсолютное дополнение}
\[
    \overline{\rho} = X \setminus \rho 
\]


\subsection{Обратное отношение}
\[
    \rho^{-1} = \{\langle x, y \rangle |~x, y \in X: \langle y, x \rangle \in \rho \}
\]


\subsection{Композиция (произведение)}
\[
    \rho_1 \circ \rho_2 = \{\langle x, y \rangle |~\exists z \in X: \langle x, z \rangle \in \rho_2, \langle z, y \rangle \in \rho_1\}
\]


\subsubsection{Пример}
\begin{gather*}
    f \circ g = f(g(x)) \\
    \langle x, z \rangle: z = g(x), \langle z, y \rangle: y = f(z) \Rightarrow \langle x, y \rangle: y = f(g(x))
\end{gather*}


\subsection{Свойства}
\begin{enumerate}
    \item $(\rho^{-1})^{-1} = \rho$
    \item $(\rho_1 \circ \rho_2)^{-1} = \rho^{-1}_2 \circ \rho^{-1}_1$
\end{enumerate}

\subsubsection{Пример}
\paragraph{Дано}
\[
    X = \{1, 2, 3, 4\}, \rho = \{\langle 1, 2 \rangle, \langle 2, 3 \rangle, \langle 2, 4 \rangle, \langle 4, 3 \rangle, \langle 4, 1 \rangle\}
\]

\paragraph{Найти}
\[
    D_{\rho}, R_{\rho}, \rho^{-1}, \rho \circ \rho^{-1}, \rho^{-1} \circ \rho, \rho^{2} = \rho \circ \rho
\]

\paragraph{Решение}
\begin{gather*}
    D_{\rho} = \{1, 2, 4\}, R_{\rho} = X \\
    \begin{tabular}{c|c|c|c|c|c}
        $\rho^{-1}$ & $\rho$ & $\rho^{-1}$ & $\rho \circ \rho^{-1}$ & $\rho^{-1} \circ \rho$ & $\rho \circ \rho$ \\
        \hline
        $\langle 2, 1 \rangle$ & $\langle 1, 2 \rangle$ & $\langle 2, 1 \rangle$ & $\langle 2, 2 \rangle$ & $\langle 1, 1 \rangle$ & $\langle 1, 3 \rangle$ \\
        $\langle 3, 2 \rangle$ & $\langle 2, 3 \rangle$ & $\langle 3, 2 \rangle$ & $\langle 3, 3 \rangle$ & $\langle 2, 2 \rangle$ & $\langle 1, 4 \rangle$ \\
        $\langle 4, 2 \rangle$ & $\langle 2, 4 \rangle$ & $\langle 4, 2 \rangle$ & $\langle 3, 4 \rangle$ & $\langle 4, 2 \rangle$ & $\langle 2, 3 \rangle$ \\
        $\langle 3, 4 \rangle$ & $\langle 4, 3 \rangle$ & $\langle 3, 4 \rangle$ & $\langle 4, 3 \rangle$ & $\langle 4, 4 \rangle$ & $\langle 2, 1 \rangle$ \\
        $\langle 1, 4 \rangle$ & $\langle 4, 1 \rangle$ & $\langle 1, 4 \rangle$ & $\langle 4, 4 \rangle$ & $\langle 2, 4 \rangle$ & $\langle 4, 2 \rangle$ \\
                               &                        &                        & $\langle 3, 1 \rangle$ &                        &                        \\
                               &                        &                        & $\langle 1, 3 \rangle$ &                        &                        \\
                               &                        &                        & $\langle 1, 1 \rangle$ &                        &       
    \end{tabular}
\end{gather*}

\subsubsection{Пример}
\paragraph{Дано}
\[
    \rho = \{\langle x, y \rangle |~x, y \in \mathbb{R}: y \divby x\} 
\]
Пример: $\langle 2, 6 \rangle \in \rho: 6 \divby 2, \langle 6, 2 \rangle \notin \rho: 2 \cancel{\divby} 6$

\paragraph{Найти}
\[
    D_{\rho}, R_{\rho}, \rho^{-1}, \rho \circ \rho^{-1}, \rho^{-1} \circ \rho, \rho^{2}
\]

\paragraph{Решение}
\begin{gather*}
    D_{\rho} = \mathbb{N} = R_{\rho} \\
    \rho^{-1} = \{\langle x, y \rangle |~x, y \in \mathbb{N}: \langle y, x \rangle \in \rho\} = x \divby y \\
    \rho \circ \rho^{-1} = 
    \{\langle x, y \rangle |~x, y \in \mathbb{N}, \exists z \in \mathbb{N}: \underbrace{\langle x, z \rangle \in \rho^{-1}}_{x \divby z}, \underbrace{\langle z, y \rangle \in \rho}_{y \divby z}\} = 
    \langle \mathbb{N}, \mathbb{N} \rangle = \mathbb{N} \times \mathbb{N} = \mathbb{N}^2 \\
    \rho^{-1} \circ \rho = 
    \{\langle x, y \rangle |~x, y \in \mathbb{N}, \exists z \in \mathbb{N}: \underbrace{\langle x, z \rangle \in \rho}_{z \divby x}, \underbrace{\langle z, y \rangle \in \rho^{-1}}_{z \divby y}\} = \mathbb{N}^2 \\
    \rho^2 = \rho \circ \rho = 
    \{\langle x, y \rangle |~x, y \in \mathbb{N}, \exists z \in \mathbb{N}: \underbrace{\langle x, z \rangle \in \rho}_{z \divby x}, \underbrace{\langle z, y \rangle \in \rho}_{y \divby z}\} = \rho \\
\end{gather*}

\subsubsection{Пример}
\paragraph{Дано}
\[
    \rho = \{\langle x, y \rangle |~x, y \in \mathbb{R}: 2x \ge 3y\}
\]

\paragraph{Найти}
\[
    D_{\rho}, R_{\rho}, \rho^{-1}, \rho \circ \rho^{-1}, \rho^{-1} \circ \rho, \rho^{2}
\]

\paragraph{Решение}
\begin{gather*}
    D_{\rho} = \mathbb{R} = R_{\rho} \\
    \rho^{-1} = \{\langle x, y \rangle |~x, y \in \mathbb{R}: \langle y, x \rangle \in \rho\} = 2y \ge 3x \\
    \rho \circ \rho^{-1} = 
    \{\langle x, y \rangle |~x, y \in \mathbb{R}, \exists z \in \mathbb{R}: \underbrace{\langle x, z \rangle \in \rho^{-1}}_{2z \ge 3x}, \underbrace{\langle z, y \rangle \in \rho}_{2z \ge 3y}\} = \mathbb{R}^2 \\
    \rho^{-1} \circ \rho = \{\langle x, y \rangle |~x, y \in \mathbb{R}, \exists z \in \mathbb{R}: \underbrace{\langle x, z \rangle \in \rho}_{2x \ge 3z}, \underbrace{\langle z, y \rangle \in \rho^{-1}}_{2y \ge 3z}\} = \mathbb{R}^2 \\
    \rho^2 = \rho \circ \rho = \{\langle x, y \rangle |~x, y \in \mathbb{R}, \exists z \in \mathbb{R}: \underbrace{\langle x, z \rangle \in \rho}_{2x \ge 3z \Rightarrow 4x \ge 6z}, \underbrace{\langle z, y \rangle \in \rho}_{2z \ge 3y \Rightarrow 6z \ge 9y}\} = 4x \ge 9y \\
\end{gather*}


\section{Основные свойства}

\subsection{Определение | Рефлексивность}
Бинарное отношение $\rho$ на множестве $X$ называется рефлексивным, если $\forall x \in X: x \rho x$.

\subsubsection{Пример}
\begin{gather*}
    \le, \mathbb{R}. ~ \forall x \in \mathbb{R}: x \le x \\
    x \rho y \Leftrightarrow x + y = 1, \mathbb{R} \\
    \langle 0, 0 \rangle \notin \rho, 0+0 \neq 1
\end{gather*}


\subsection{Определение | Симметричность}
Бинарное отношение $\rho$ на множестве $X$ называется симметричным, если $\forall x, y \in X: x \rho y \Rightarrow y \rho x$.

\subsubsection{Пример}
\begin{gather*}
    \le, \mathbb{R}. ~ \forall x, y \in X: \\
    x \rho y \Leftrightarrow x + y = 1, \mathbb{R} \\
    \le, \mathbb{R}. ~ x + y = 1 \Rightarrow y + x = 1 \\
    1 \le 2 \nRightarrow 2 \le 1
\end{gather*}


\subsection{Определение | Антисимметричность}
Бинарное отношение $\rho$ на множестве $X$ называется антисимметричным, если 
\[
    \forall x, y \in X: x \rho y, y \rho x \Rightarrow x = y.
\]

\subsubsection{Пример}
\begin{gather*}
    \le, \mathbb{R}. ~ x \le y, y \le x \Leftrightarrow x = y \\
    x \rho y \Leftrightarrow x + y = 1, \mathbb{R} \\
    \langle 0, 1 \rangle \in \rho, \langle 1, 0 \rangle \in \rho, 1 \neq 0
\end{gather*}


\subsection{Определение | Асимметричность}
Бинарное отношение $\rho$ на множестве $X$ называется асимметричным, если
\[
    \forall x, y \in X: x \rho y \Rightarrow y \overline{\rho} x.
\]

\subsubsection{Пример}
\begin{gather*}
    <, \mathbb{R}. ~ x < y ~\Rightarrow y \cancel{<} x \\
    x \rho y \Leftrightarrow x + y = 1, \mathbb{R} \\
    \langle 0, 1 \rangle, \langle 1, 0 \rangle \in \rho
\end{gather*}


\subsection{Определение | Транзитивность}
Бинарное отношение $\rho$ на множестве $X$ называется транзитивным, если
\[
\forall x, y, z \in X: x \rho y, y \rho z \Rightarrow x \rho z.
\]

\subsubsection{Пример}
\begin{gather*}
    <, \mathbb{R}. ~ x < y, y < z \Rightarrow x < z \\
    \text{любовь, люди} \\
    \text{Петя любит Машу, Маша любит Васю } \cancel{\Rightarrow} \text{ Петя любит Васю}
\end{gather*}

\subsection{Практика} 

\subsubsection{Пример | Расчётно-графическая работа: задание 2}
\paragraph{Дано}
\[
    X = \{1, 2, 3, 4\}, \rho = \{\langle 1, 1 \rangle, \langle 1, 2 \rangle, \langle 2, 1 \rangle, \langle 2, 2 \rangle, \langle 2, 3 \rangle, \langle 1, 3 \rangle, \langle 3, 4 \rangle, \langle 4, 4 \rangle\} \\
\]

\paragraph{Найти}
\[
    D_{\rho}, R_{\rho}, \rho^{-1}, \rho^2; \text{Р, С, Ант, Ас, Т --- ?}
\]

\paragraph{Решение}
\begin{gather*}
    D_{\rho} = X = R_{\rho} \\
    \begin{tabular}{c|c|c}
        $\rho$ & $\rho^{-1}$ & $\rho^2$ \\
        \hline
        $\langle 1, 1 \rangle$ & $\langle 1, 1 \rangle$ & $\langle 1, 1 \rangle$ \\
        $\langle 1, 2 \rangle$ & $\langle 2, 1 \rangle$ & $\langle 1, 2 \rangle$ \\
        $\langle 2, 1 \rangle$ & $\langle 1, 2 \rangle$ & $\langle 1, 3 \rangle$ \\
        $\langle 2, 2 \rangle$ & $\langle 2, 2 \rangle$ & $\langle 2, 1 \rangle$ \\
        $\langle 2, 3 \rangle$ & $\langle 3, 2 \rangle$ & $\langle 2, 2 \rangle$ \\
        $\langle 1, 3 \rangle$ & $\langle 3, 1 \rangle$ & $\langle 2, 3 \rangle$ \\
        $\langle 3, 4 \rangle$ & $\langle 4, 3 \rangle$ & $\langle 2, 4 \rangle$ \\
        $\langle 4, 4 \rangle$ & $\langle 4, 4 \rangle$ & $\langle 1, 4 \rangle$ \\
                               &                        & $\langle 3, 4 \rangle$ \\
                               &                        & $\langle 4, 4 \rangle$
    \end{tabular}
    \\
\end{gather*}
Рефлексивность --- нет, так как $\langle 3, 3 \rangle \notin \rho$ \\
Симметричность --- нет, так как $\langle 2, 3 \rangle \in \rho, \langle 3, 2 \rangle \notin \rho$ \\
Антисимметричность --- нет, так как $\langle 1, 2 \rangle \in \rho, \langle 2, 1 \rangle \in \rho, 2 \neq 1$ \\
Асимметричность --- нет, так как $\langle 1, 1 \rangle \in \rho$ \\
Транзитивность --- нет, так как $\langle 2, 3 \rangle \in \rho, \langle 3, 4 \rangle \in \rho, \langle 2, 4 \rangle \notin \rho$ \\
\includegraphics[scale=0.3]{images/part03_03.jpg}



\section{Другие свойства}

\subsection{Определение | Эквивалентность}
Бинарное отношение $\rho$ на множестве $X$ называется отношением эквивалентности, \\
если $\rho$ --- рефлексивно, симметрично, транзитивно.


\subsection{Определение | Частичный порядок}
Бинарное отношение $\rho$ на множестве $X$ называется отношением частичного порядка, \\
если $\rho$ --- рефлексивно, антисимметрично, транзитивно.


\subsection{Определение | Квазипорядок}
Бинарное отношение $\rho$ на множестве $X$ называется отношением квазипорядка, \\
если $\rho$ --- рефлексивно, транзитивно.


\subsection{Определение | Строгий порядок}
Бинарное отношение $\rho$ на множестве $X$ называется отношением строгого порядка, \\
если $\rho$ --- ассимметрично, транзитивно.


\subsection{Определение | Линейный порядок}
Бинарное отношение $\rho$ на множестве $X$ называется отношением линейного порядка, \\
если $\rho$ --- частичный порядок, у которого все элементы сравнимы.
\[
\forall x, y: x \rho y \text{ или } y \rho x.
\]



\section{Условия}

\subsection{Рефлексивность}
\[
    d = \{\langle x, x \rangle |~\forall x \in X\}, \boxed{d \subseteq \rho}
\]
$d$ --- отношение-диагональ \\
\includegraphics[scale=0.3]{images/part03_04.png}

\subsubsection{Пример}
\[
    d = \{\langle 1, 1 \rangle, \langle 2, 2 \rangle, \langle 3, 3 \rangle, \langle 4, 4 \rangle\}
\]


\subsection{Симметричность}
\[
    \boxed{\rho = \rho^{-1}}
\]
\includegraphics[scale=0.3]{images/part03_05.png}


\subsection{Антисимметричность}
\[
    \boxed{\rho \cap \rho^{-1} \subseteq d}
\]
\includegraphics[scale=0.3]{images/part03_06.png}

\subsubsection{Пример}
\[
    \rho \cap \rho^{-1} = \{\langle 1, 1 \rangle, \langle 1, 2 \rangle, \langle 2, 1 \rangle, \langle 2, 2 \rangle, \langle 4, 4 \rangle\}
\]


\subsection{Асимметричность}
\[
    \boxed{\rho \cap \rho^{-1} = \varnothing}
\]
\includegraphics[scale=0.3]{images/part03_07.png}


\subsection{Транзитивность}
\begin{gather*}
    x \rho y, y \rho z \Rightarrow x \rho z \\
    x \rho y, y \rho z \Rightarrow x \rho^2 z \\
    \boxed{\rho^2 \subseteq \rho} \\
    \rho = \{\langle 1, 2 \rangle\} \text{ --- транзитивно} \\
    \rho^2 = \varnothing
\end{gather*}



\subsection{Практика}
\paragraph{Дано}
\begin{enumerate}
    \item $\rho_1 \cap \rho_2$
    \item $\rho_1 \cup \rho_2$
\end{enumerate}

\subsubsection{Пример}
\paragraph{Узнать}
Рефлексивны на $X$?

\paragraph{Решение для 1 и 2}
\[
    \forall x \in X, \langle x, x \rangle \in \rho_1, \rho_2 \underset{\text{рефл.}}{\Rightarrow} \langle x, x \rangle \in \rho_1 \cap \rho_2, \rho_1 \cup \rho_2
\]


\subsubsection{Пример}
\paragraph{Узнать}
Симметричны на $X$?

\paragraph{Решение для 1}
\[
    \langle x, y \rangle \in \rho_1 \cap \rho_2 \Rightarrow \langle x, y \rangle \in \rho_1, \rho_2 \underset{\text{асимм.}}{\Rightarrow} \langle y, x \rangle \in \rho_1, \rho_2 \Rightarrow \langle y, x \rangle \in \rho_1 \cap \rho_2
\]

\paragraph{Решение для 2}
\[
    \langle x, y \rangle \in \rho_1 \cup \rho_2 \Rightarrow \langle x, y \rangle \in \rho_1 \cup \rho_2
\]


\subsubsection{Пример}
\paragraph{Узнать}
Антисимметричны на $X$?

\paragraph{Решение для 1}
\begin{gather*}
    \text{Если } x \neq y: \langle x, y \rangle \in \rho_1 \cap \rho_2 \Rightarrow 
    \langle x, y \rangle \in \rho_1, \rho_2 \underset{\text{ант.}}{\Rightarrow} 
    \langle y, x \rangle \notin \rho_1, \rho_2 \Rightarrow \langle y, x \rangle \notin \rho_1 \cap \rho_2
\end{gather*}

\paragraph{Решение для 2}
Нет, так как
\[
    \rho_1 = \{\langle 1, 2 \rangle\}, \rho_2 = \{\langle 2, 1 \rangle\} \text{ --- ант.}, \rho_1 \cup \rho_2 = \{\langle 1, 2 \rangle, \langle 2, 1 \rangle\} \text{ --- не ант.}
\]


\subsubsection{Пример}
\paragraph{Узнать}
Асимметричны на $X$?

\paragraph{Решение для 1 и 2}
Аналогично прошлому премеру


\subsubsection{Пример}
\paragraph{Узнать}
Транзитивны на $X$?

\paragraph{Решение для 1}
\begin{gather*}
    \langle x, y \rangle \in \rho_1 \cap \rho_2, \langle y, z \rangle \in \rho_1 \cap \rho_2 \Rightarrow \langle x, y \rangle \in \rho_1, \rho_2; \langle y, z \rangle \in \rho_1, \rho_2 \Rightarrow \\
    \underset{\text{тр. }\rho_2}{\Rightarrow} \langle x, z \rangle \in \rho_1, \rho_2 \Rightarrow \langle x, z \rangle \in \rho_1 \cap \rho_2
\end{gather*}


\paragraph{Решение для 2}
\[
    \rho_1 = \{\langle 1, 2 \rangle\}, \rho_2 = \{\langle 2, 3 \rangle\} \text{ --- тр.}, \rho_1 \cup \rho_2 = \{\langle 1, 2 \rangle, \langle 2, 3 \rangle\} \text{ --- нет}
\]




\newpage
\paragraph{2024.10.29 (3, 4)}
\part{Специальные типы бинарных отношений}

\section{Отношения частичного порядка}

\subsection{Определение}
$\rho$ на $X$ --- частичный порядок, если $\rho$ --- рефлексивно, антисимметрично, транзитивно.


\subsection{Определение}
Частичный порядок, у которого все элементы сравнимы, называется линейным порядком (ЛП).
\[
    \forall a, b \in X: a \rho b \text{ или } b \rho a
\]
Обозначение --- $\preccurlyeq (\succcurlyeq)$

\subsubsection{Пример}
\[
    =, \mathbb{R}
\]
\begin{gather*}
    \forall x \in \mathbb{R}: x = x \\
    \forall x, y \in \mathbb{R}: x = y, y = x \Rightarrow x = y \\
    \forall x, y, z \in \mathbb{R}: x = y, y = z \Rightarrow x = z
\end{gather*}

\subsubsection{Пример}
\[
    \le, \mathbb{R}
\]
\begin{gather*}
    \forall x \in \mathbb{R}: x \le x \\
    \forall x, y \in \mathbb{R}: x \le y, y \le x \Rightarrow x = y \\
    \forall x, y, z \in \mathbb{R}: x \le y, y \le z \Rightarrow x \le z
\end{gather*}

\subsubsection{Пример}
\[
    \subseteq, \mathbb{U}
\]
\begin{gather*}
    A \subseteq A, A \subseteq B, B \subseteq A \Rightarrow A = B \\
    A \subseteq B, B \subseteq C \Rightarrow A \subseteq C
\end{gather*}
\includegraphics[scale=0.4]{images/part04_01.png}

\paragraph{Примечание}
~\\
\begin{gather*}
    \forall A_i, A_j \text{ --- сравнимы}, где A_i \subseteq A_j \text{ или } A_j \subseteq A_i \\
    \forall B_i, B_j \text{ --- сравнимы}, где B_i \subseteq B_j \text{ или } B_j \subseteq B_i \\
    \forall A_i, B_i \text{ --- не сравнимы}
\end{gather*}
(фото 1, 2, 3)



\section{Отношение Парето (подраздел частичного порядка)}

\subsection{Определение}
\begin{gather*}
    \langle x, y \rangle \text{ П } \langle a, b \rangle \Leftrightarrow (x \ge a, y > b) \text{ или } (x > a, y \ge b)    
\end{gather*}
 

\subsection{Диаграмма Хассе}
\begin{gather*}
    x \preccurlyeq y \\
    x \prec y \Leftrightarrow x \preccurlyeq y, x \neq y \\
    y \text{ покрывает } x \text{(непосредственно предшествует)}:
    \nexists u \in X: x \prec u \prec y
\end{gather*}
\includegraphics[scale=0.2]{images/part04_02.png}
\begin{gather*}
    x \preccurlyeq y \Leftrightarrow x = u_1 \prec u_2 \prec \dots \prec u_n = y, \text{где } u_{i + 1} \text{ покрывает } u_i
\end{gather*}



\subsubsection{Пример}
\[
    X = \{ 1, 2, 3, 5, 6, 10, 15, 30 \}~~~~~ x \rho y \Leftrightarrow y \divby x
\]
\includegraphics[scale=0.2]{images/part04_03.png}


\subsection{Определение}
\begin{gather*}
    x \in \min_{X} \rho \text{, если } \nexists y: y \prec x \\
    x \in \max_{X} \rho \text{, если } \nexists y: y \succ x \\
    x \text{ --- наименьший}: \min \text{, сравнимый со всеми;} \\
    x \text{ --- наибольший}: \max \text{, сравнимый со всеми,} \\
    \text{то есть выполняется условие } \forall y: x \rho y
\end{gather*}

\subsubsection{Примечение}
Наименьший (наибольший) либо $\nexists$, либо единственный.

\subsubsection{Пример}
\includegraphics[scale=0.3]{images/part04_04.png} \\
\begin{itemize}
    \item Наименьших нет, так как $1$ и $2$ не сравнимы.
    \item Единственность: предположим, $\exists a \neq b$ --- наименьший
\end{itemize}
\begin{align*}
\begin{rcases}
    a \text{ --- наименьший } \Rightarrow a \rho b \\
    b \text{ --- наименьший } \Rightarrow b \rho a
\end{rcases}
\Rightarrow a = b
\end{align*}

\subsubsection{Пример}
\includegraphics[scale=0.2]{images/part04_05.png}
\begin{gather*}
    \rho = \{ \langle a, a \rangle, \langle b, b \rangle, \langle c, c \rangle, \langle d, d \rangle, \langle e, e \rangle, \langle f, f \rangle, \langle g, g \rangle, \\
              \langle a, c \rangle, \langle a, e \rangle, \langle a, f \rangle, \langle a, g \rangle, \langle b, f \rangle, \langle b, g \rangle, \\ 
              \langle d, c \rangle, \langle d, e \rangle, \langle d, f \rangle, \langle d, g \rangle, \langle c, f \rangle, \langle c, g \rangle, \langle e, f \rangle, \langle e, g \rangle, \langle f, g \rangle \} \\
    \min_X \rho = {a, b, d} \\
    \max_X \rho = {g} \\
    \text{Наименьший: —} \\
    \text{Наибольший: } g
\end{gather*}

\subsubsection{Пример}
\includegraphics[scale=0.2]{images/part04_06.png}
\begin{gather*}
    \rho = \{ \langle a, a \rangle, \langle b, b \rangle, \langle c, c \rangle, \langle d, d \rangle, \langle e, e \rangle, \langle f, f \rangle, \langle g, g \rangle, \textcolor{blue}{\langle k, k \rangle}, \\
              \langle a, c \rangle, \langle a, e \rangle, \langle a, f \rangle, \langle a, g \rangle, \langle b, f \rangle, \langle b, g \rangle, \\ 
              \langle d, c \rangle, \langle d, e \rangle, \langle d, f \rangle, \langle d, g \rangle, \langle c, f \rangle, \langle c, g \rangle, \langle e, f \rangle, \langle e, g \rangle, \langle f, g \rangle \} \\
    \min_X \rho = {a, b, d, \textcolor{blue}{k}} \\
    \max_X \rho = {g, \textcolor{blue}{k}} \\
    \text{Наименьший: —} \\
    \text{Наибольший: } \textcolor{blue}{\text{—}}
\end{gather*}


\subsection{Алгоритм достроения до линейного порядка}
\begin{align*}
    X_1 &= \min_X \rho \text{ --- произвольно упорядоченное} \\
    X_2 &= \min_{X \setminus X_1} \rho \text{ --- произвольно упорядоченное} \\
    \vdots \\
    X_k &= \min_{X \setminus (X_1 \cup \dots \cup X_{k-1})} \rho \text{ --- произвольно упорядоченное} \\
    X_{k + 1} &= \varnothing
\end{align*}

\subsubsection{Пример}
\begin{align*}
    X_1 &= \min_X \rho = \{a, b, d\} \\
    X_2 &= \min_{X \setminus X_1} \rho = \{c, e\} \\
    X_3 &= \min_{X \setminus (X_1 \cup X_2)} \rho = \{f\} \\
    X_4 &= \min_{X \setminus (x_1 \cup X_2 \cup X_3)} \rho = \{g\}
\end{align*}
\includegraphics[scale=0.2]{images/part04_07.png}



\section{Отношение эквивалентности}
$\rho$ на $X$ --- EQ, если $\rho$ --- рефлексивно, симметрично, транзитивно. 

\subsubsection{Пример}
\[
    =, \mathbb{R}
\]
\begin{gather*}
    \forall x \in \mathbb{R}: x = x \\
    \forall x, y \in \mathbb{R}: x = y \Rightarrow y = x \\
    \forall x, y, z \in \mathbb{R}: x = y, y = z \Rightarrow x = z
\end{gather*}

\subsubsection{Пример}
Принадлежность к студенческой группе
\[
    x \rho y \Leftrightarrow x, y \text{ учатся в одной группе}
\]


\subsection{Определение}
Класс эквивалентности, порождаемый элементом $x \in X$.
\begin{gather*}
    [x_p] = [x] \text{ по отношению эквивалентности } \rho \text{ на } X: \\
    [x] = \{y \in X~|~x\rho y\}.
\end{gather*}

\subsubsection{Пример}
\[
    [x]_= =\{x\}
\]

\subsubsection{Пример}
Студенческая группа
$[$Иванов$]$ = $\{$Иванов, Петров, \dots$\} = [$Петров$]$


\subsection{Определение}
\[
    x \rho y \Leftrightarrow x \underbrace{\equiv} y(\text{mod } n), n \in \mathbb{N}; N \cup \{0\} = N_0
\]
сравнимость по модулю $n$ --- один. остаток от деления на $n$

\subsubsection{Пример}
\begin{gather*}
    x \rho y \Leftrightarrow x \equiv y (\text{mod } 3), N_0. \\
    \rho = \{ <0, 3>, <21, 15>, \dots, <1, 4>, <13, 7>, \dots\, <2, 5>, <17, 11>, \dots\} \\
    [0]_{\equiv 3} = \{0, 3, 6, 9, \dots\} = [3] = [6] = \dots \\
    [1] = \{1, 4, 7, 10, \dots\} = [4] = \dots \\
    [2] = \{2, 5, 8, 11, \dots\} = [8] = \dots \\
    [3] = [0] \\
    \dots \\
    N_0 = [0] \cup [1] \cup [2] \\
    [i] \cap [j] = \varnothing \\
    i \neq j ~ \{0, 1, 2\}
\end{gather*}

\subsection{Определение}
Фактор-множество --- совокупность (подмножеств разбиения) классов эквивалентности по $\rho$ --- EQ на $X$ --- $X/\rho$.

\subsubsection{Пример}
\[
    N_0/\equiv_3 \{[0], [1], [2]\}
\]



\section{Практика}

\subsection{Пример | Расчётно-графическая работа: задание 3}
\begin{align*}
    &x R_f y \Leftrightarrow f(x) = f(y) \\
    &X = [0; 5] 
\end{align*}
\includegraphics[scale=0.2]{images/part04_08.png}
\begin{align*}
    &\forall x \in (0; 1): [x] = \{x, 2 - x, 3 + x\} \\
    &\forall x \in (1; 2): [x] = \{2 - x, x, 5 - x\} \\
    &\forall x \in [2; 3] \cup \{0\}: [x] = \{0\} \cup [2, 3] \\
    &\forall x \in (3; 4): [x] = \{x - 3, 5 - x, x\} \\
    &\forall x \in (4; 5]: [x] = \{x\} \\
    &[1] = \{1;4\} = [4]
\end{align*}


\subsection{Пример | Расчётно-графическая работа: задание 3}
\begin{align*}
    &x R_f y \Leftrightarrow f(x) = f(y) \\
    &X = [0; 5] 
\end{align*}
\includegraphics[scale=0.2]{images/part04_09.png}
\begin{align*}
    &\forall x \in (0; 1]: [x] = \{x, 2 - x, 2 + x, 4 - x\} \\
    &\forall x \in (1; 2): [x] = \{x, 2 - x, 2 + x, 4 - x\} \\
    &\forall x \in (2; 3]: [x] = \{x, 4 - x, 6 - x, x - 2\} \\
    &\forall x \in (3; 4): [x] = \{x, 4 - x, 6 - x, x - 2\} \\
    &\forall x \in \{0; 2\} \cup (4; 5]: [x] = \{([4; 5] \cup 0 \cup 2)\} \\
\end{align*}



\section{Функция}

\subsection{Определение}
Бинарное отношение называется функцией, если
\begin{gather*}
    \langle x, y_1 \rangle \in f, \langle x, y_2 \rangle \in f \Rightarrow y_1 = y_2 \\
    \langle x, y \rangle \in f, x f y, y = f(x)
\end{gather*}

\subsubsection{Пример}
\begin{gather*}
    f = \{\langle 1, 2 \rangle, \langle 1, 3 \rangle\} \text{ --- не функция} \\
    f = \{\langle 1, 2 \rangle, \langle 2, 3 \rangle\} \text{ --- функция}
\end{gather*}
\includegraphics[scale=0.2]{images/part04_10.png}

\subsubsection{Пример}
\[
    f = \{\langle x, y \rangle ~|~ \mathbb{R}, y = \sqrt{x - 1}\} ~~~ D_f = [1; +\infty) ~~~ R_f = [0; +\infty)
\]

\subsection{Определение}
\includegraphics[scale=0.2]{images/part04_11.png}


\subsection{Свойства}

\subsubsection{Сюръекция}
\[
    f \text{ --- сюръективна, если } \forall y \in Y ~ \exists x \in X: \langle x, y \rangle \in f.
\]

\subsubsection{Инъекция}
\[
    f \text{ --- инъъективна, если } \langle x_1, y \rangle \in f, \langle x_2, y \rangle \in f \Rightarrow x_1 = x_2.
\]

\subsubsection{Биекция}
\[
    f \text{ --- биективна, если } f \text{ --- сюръективна и инъективна.}
\]

\subsubsection{Примеры для всех свойств}
\[
    f: \mathbb{R} \rightarrow \mathbb{R}
\]
\includegraphics[scale=0.2]{images/part04_12.png}

\subsubsection{Пример}
\[
    f: X \rightarrow Y
\]
\paragraph{Пункт 1}
\begin{gather*}
    f^{-1} \text{ --- функция?} \\
    f = \{\langle 1, 3 \rangle, \langle 2, 3 \rangle\} \text{ --- функция} \\
    f^{-1} = \{\langle 3, 1 \rangle, \langle 3, 2 \rangle\} \text{ --- не функция} \\
    f^{-1}: Y \rightarrow X \text{ --- функция } \Leftrightarrow f \text{ --- биекция}
\end{gather*}
\paragraph{Пункт 2}
\begin{gather*}
    f: X \rightarrow Y, g: Y \rightarrow Z \text{ --- функция} \Rightarrow g \circ f: X \rightarrow Z \text{ --- функция}
\end{gather*}



\subsubsection{Пример}
\[
    f(A \cap B) \neq_{\subseteq} f(A) \cap f(B)
\]
\paragraph{Пункт 1}
\[
    f(A \cap B) \subseteq f(A) \cap f(B)
\]
\begin{gather*}
    \text{Пусть } y \in f(A \cap B) \Rightarrow 
    \exists x \in A \cap B \Rightarrow 
    x \in A \text{ и } x \in B \Rightarrow
    y \in f(A) \text{ и } y \in f(B) \Rightarrow
    y \in f(A) \cap f(B).
\end{gather*}
\paragraph{Пункт 2}
\[
    f(A) \cap f(B) \subseteq f(A \cap B)
\]
\begin{gather*}
    \text{Пусть } y \in f(A) \cap f(B) \Rightarrow
    y \in f(A) \text{ и } y \in f(B) \nRightarrow
    \exists x \in A \text{ и } x \in B \Rightarrow
    x \in A \cap B \Rightarrow y \in f(A \cap B).
\end{gather*}




\newpage
\paragraph{2024.11.12 (3)}
\part{Мощность множеств}

\section{Введение}

\subsection{Определение}
Два множества А и В называются равномощными, если существует биекция, отображающая одно множество на другом:
\[
    A \sim B \text{ --- равномощны, если } \exists \text{ биекция } \varphi: A \rightarrow B 
\]
Равномощность --- отношение эквивалентности

\subsection{Свойства}
\begin{itemize}
    \item Рефлексивно, если $\varphi: A \rightarrow A$ ~~~ $\rho_A$ --- биекция 
    \item Симметрично, если $\varphi$ --- биекция ~~~ $\varphi: A \rightarrow B \Rightarrow \varphi^{-1}: B \rightarrow A$ --- биекция
    \item Транзитивно, если $\varphi$, $\psi$ --- биекции $\Rightarrow \varphi \circ \psi$ --- биекция
\end{itemize}



\section{Типы}

\subsection{Конечные множества}
\begin{gather*}
    A = \{a_1, \dots, a_k\}, B = \{b_1, \dots, b_m\} \\
    A \sim B \Rightarrow k = m \\
\end{gather*}
\begin{itemize}
    \item $k > m$ --- не инъективно
    \item $m > k$ --- не сюръективно
\end{itemize}


\subsection{Счётные множества}
\begin{gather*}
    A \sim \mathbb{N}, \mathbb{Z} \sim \mathbb{N} \\
    \mathbb{Z}: 0, 1, -1, 2, -2, \dots ~~~ z \le 0 \rightarrow n = 2|z| + 1 \\
    \mathbb{N}: 1, 2,  3, 4,  5, \dots ~~~ z  >  0 \rightarrow n = 2z
\end{gather*}
\includegraphics[scale=0.3]{images/part05_01.png}
\includegraphics[scale=0.3]{images/part05_02.png}


\subsection{Континуум}
\includegraphics[scale=0.3]{images/part05_03.png} \\
\includegraphics[scale=0.3]{images/part05_04.png}



\section{Практика}

\subsection{Пример | И. А. Лавров, Л. Л. Максимова: часть 1, параграф 4, номер 20}
\paragraph{Доказать}
Любое множество попарно не пересекающихся открытых интервалов на множестве действительных чисел не более чем счётно.

\paragraph{Доказательство}
~\\
\includegraphics[scale=0.3]{images/part05_05.png}


\subsection{Пример | И. А. Лавров, Л. Л. Максимова: часть 1, параграф 4, номер 23}
\paragraph{Доказать}
Множество точек разрыва монотонной функции на множестве действительных чисел не более чем счётно.

\paragraph{Доказательство}
~\\
\includegraphics[scale=0.3]{images/part05_06.png}


\section{Теорема Кантора-Бернштейна}
\subsection{Случай 1}
\[
    A, B ~~~ A_1 \subseteq A, B_1 \subseteq B: A_1 \sim B, B_1 \sim A \Rightarrow A \sim B
\]

\subsection{Случай 2}
\[
    A_1 \sim B, \nexists B_1 \sim A \Rightarrow m(A) > m(B)
\]

\subsection{Случай 3}
\[
    B_1 \sim A, \nexists A_1 \sim B \Rightarrow m(B) > m(A)
\]

\subsection{Случай 4}
\[
    \nexists A_1 \sim B, \nexists B_1 \sim A \Rightarrow A, B \text{ --- не сравнимы?}
\]

\section{Лемма Цермело}
Любые множества сравнимы по мощности.
\begin{gather*}
    \forall A: |A| \ge 2 ~ P(A) = 2^A \text{ --- множество-степень} \\
    m(P(A)) > m(A) \\
    m(k) < m(\mathbb{N}) < m(\mathbb{R}) < \dots
\end{gather*}


\subsection{Пример | И. А. Лавров, Л. Л. Максимова: часть 1, параграф 4, номер 25}
\includegraphics[scale=0.3]{images/part05_07.png}




\newpage
\part{Математическая логика \\ Логика высказываний}

\section{Введение}

\subsection{Определение}
Высказывание --- предложение, о котором есть смысл говорить истинно оно или ложно.



\section{Операции}

\subsection{Отрицание (не)}
\begin{tabular}{c|c}
    $P$ & $\neg P$ \\
    \hline
    И & Л \\
    Л & И
\end{tabular}


\subsection{Конъинкция (и)}
\begin{tabular}{c|c|c}
    $P$ & $Q$ & $P ~\&~ Q$ \\
    \hline
    Л & Л & Л \\
    Л & И & Л \\
    И & Л & Л \\
    И & И & И
\end{tabular}


\subsection{Дизъюнкция (или)}
\begin{tabular}{c|c|c}
    $P$ & $Q$ & $P \vee Q$ \\
    \hline
    Л & Л & Л \\
    Л & И & И \\
    И & Л & И \\
    И & И & И
\end{tabular}


\subsection{Импликация}
\begin{itemize}
    \item если \dots, то \dots
    \item следует
    \item влечёт
\end{itemize}
\begin{tabular}{c|c|c}
    $P$ & $Q$ & $P \supset Q$ \\
    \hline
    Л & Л & И \\
    Л & И & И \\
    И & Л & Л \\
    И & И & И
\end{tabular}

\subsubsection{Пример}
\begin{itemize}
    \item $P$ --- 'Пацан сказал'
    \item $Q$ --- 'Пацан сделал'
\end{itemize}


\subsection{Эквивалентация (тогда и только тогда)}
\begin{tabular}{c|c|c}
    $P$ & $Q$ & $P \sim Q$ \\
    \hline
    Л & Л & И \\
    Л & И & Л \\
    И & Л & Л \\
    И & И & И
\end{tabular}


\subsection{Исключающие 'или'}
\begin{itemize}
    \item или \dots, или \dots
    \item либо \dots, либо \dots
\end{itemize}

\begin{tabular}{c|c|c}
    $P$ & $Q$ & $P + Q$ \\
    \hline
    Л & Л & Л \\
    Л & И & И \\
    И & Л & И \\
    И & И & Л
\end{tabular}


\section{Пример}
\[
    \underbrace{\underbrace{\neg \underbrace{(\underbrace{\neg x}_{(3)} \supset \underbrace{(y + \underbrace{\neg z}_{(1)})}_{(2)})}_{(4)}}_{(5)} \sim \underbrace{(\underbrace{\neg y}_{(7)} \vee \underbrace{(z ~\&~ x)}_{(6)})}_{(8)}}_{(9)}
\]

\begin{tabular}{c|c|c|c|c|c|c|c|c|c|c|c}
    $x$ & $y$ & $z$ & (1) & (2) & (3) & (4) & (5) & (6) & (7) & (8) & (9) \\
    \hline
    И & И & И & Л & И & Л & И & Л & И & Л & И & Л \\
    И & И & Л & И & Л & Л & И & Л & Л & Л & Л & И \\
    И & Л & И & Л & Л & Л & И & Л & И & И & И & Л \\
    И & Л & Л & И & И & Л & И & Л & Л & И & И & Л \\
    Л & И & И & Л & И & И & И & Л & Л & Л & Л & И \\
    Л & И & Л & И & Л & И & Л & И & Л & Л & Л & Л \\
    Л & Л & И & Л & Л & И & Л & И & Л & И & И & И \\
    Л & Л & Л & И & И & И & И & Л & Л & И & И & Л
\end{tabular}



\newpage
\paragraph{2024.11.12 (4)}
\part{Математическая логика \\ Тавтологии. Правильные рассуждения}

\section{Введение}
\[
A <x_1, \dots x_n>
\]

\subsection{Определение}
Формула называется тождественно истинной (тавтологией), если она принимает значение истинна на каждой оценке списка. 


\subsection{Определение}
Формула называется выполнимой, если она принимает значение истинна хотя бы на одной оценке списка.


\subsection{Определение}
Формула называется опровержимой, если она принимает значение ложь хотя бы на оценке списка.


\subsection{Определение}
Формула называется тождественно ложной, если она принимает значение ложь на каждой оценке списка.


\section{Тавтологии}
\begin{itemize}
    \item $A \vee \neg A$
    \item $A \supset A$
    \item $(A \supset B) \supset ((B \supset C) \supset (A \supset C))$ --- ценное рассуждение \\
    \paragraph{Доказательство}
    ~\\
    \includegraphics[scale=0.3]{images/part07_01.png}
    \item $(A \supset (B \supset C)) \supset ((A \supset B) \supset (A \supset C))$ \\
    \paragraph{Доказательство}
    ~\\
    \includegraphics[scale=0.3]{images/part07_02.png}
\end{itemize}



\section{Посылки, рассуждения, заключения}

\subsection{Определения}
\begin{itemize}
    \item $P_1, \dots, P_n$ --- посылки
    \item $D$ --- заключение
    \item Рассуждение --- правильное, если $(P_1 ~ ~\&~ ~ P_2 ~ ~\&~ ~ \dots ~ ~\&~ ~ P_n) \supset D$ --- тождественно инстинно
    \item Схема --- $\frac{P_1, \dots, P_n}{D}$
\end{itemize}


\subsubsection{Примеры}
\paragraph{Дано}
~\\
Проверить правильность следующего рассуждения. \\
Если Джонс не встречал этой ночью Смита, то либо Смит был убийцой, либо Джонс лжёт. \\
Если Смит не был убийцей, то Джонс не встречал Смита этой ночью, и убийство имело место после полуночи. \\
Если убийство было совершено после полуночи, то либо Смит был убийцей, либо Джонс лжёт. \\
Следовательно, Смит был убийцей.

\paragraph{Высказывания}
\begin{itemize}
    \item $A$ --- 'Джонс не встречал ночью Смита'
    \item $B$ --- 'Смит --- убийца'
    \item $C$ --- 'Джонс лжёт'
    \item $D$ --- 'Убийство после полуночи'
\end{itemize}

\paragraph{Посылки}
\begin{itemize}
    \item $P_1 = A \supset (B + C)$
    \item $P_2 = \neg B \supset (A ~\&~ D)$
    \item $P_3 = D \supset (B + C)$
\end{itemize}

\paragraph{Заключение}
~\\
$Q = B$

\paragraph{Решение}
~\\
\includegraphics[scale=0.3]{images/part07_03.png}

\paragraph{Ответ}
~\\
Рассуждение правильно


\subsection{Пример}
\paragraph{Дано}
~\\
Проверить правильность следующего рассуждения. \\
Если Иван приходит домой не поздно и в хорошем настроении, то он смотрит телевизор. \\
Когда он смотрит телевизор, то всегда что-нибудь ест. \\
Иван пришёл домой не поздно и ничего не ест. \\
Следовательно, он пришёл домой в плохом настроении.

\paragraph{Высказывания}
\begin{itemize}
    \item $A$ --- 'Иван пришёл не поздно'
    \item $B$ --- 'Иван пришёл в хорошем настроении'
    \item $C$ --- 'Иван смотрит телевизор'
    \item $D$ --- 'Иван что-нибудь ест'
\end{itemize}

\paragraph{Посылки}
\begin{itemize}
    \item $P_1 = (A ~\&~ B) \supset C$
    \item $P_2 = C \supset D$
    \item $P_3 = A ~\&~ \neg D$
\end{itemize}

\paragraph{Заключение}
~\\
$Q = \neg B$

\paragraph{Решение}
~\\
\includegraphics[scale=0.3]{images/part07_04.png}

\paragraph{Ответ}
~\\
Рассуждение правильное




\newpage
\part{Нормальные формы формул}

\section{Введение}

\subsection{Определение}
Символы конъюнкции и дизъюнкции называются двойственными
\begin{gather*}
    ~\&~, \vee \text{ --- двойственные символы} \\
    A^* \text{ --- двойственна } A, \text{ если в } A: ~\&~ \leftrightarrow \vee
\end{gather*}

\subsubsection{Пример}
\[
    (\neg x ~\&~ (y \vee z))^* \equiv \neg x \vee (y ~\&~ \neg z)
\]

\subsection{Теорема (пр. двойств.)}
\begin{itemize}
    \item $A \equiv B \Rightarrow A^* \equiv B^*$
    \item $A \supset B \equiv \neg A \vee B \equiv \neg (A ~\&~ \neg B)$
    \item $A \sim B \equiv (A ~\&~ B) \vee (\neg A ~\&~ \neg B) \equiv (A \supset B) ~\&~ (B \supset A) \equiv (\neg A \vee B) ~\&~ (A \vee \neg B)$
    \item $A + B \equiv \neg (A \sim B) \equiv (A ~\&~ \neg B) \vee (\neg A ~\&~ B) \equiv (A \vee B) ~\&~ (\neg A \vee \neg B)$
\end{itemize}



\section{Свойства}

\subsection{Коммутативность}
\[
    A \vee B \equiv B \vee A
\]

\subsection{Ассоциативность}
\[
    A \vee (B \vee C) \equiv (A \vee B) \vee C 
\]

\subsection{Дистрибутивность}
\[
    A \vee (B ~\&~ C) \equiv (A \vee B) ~\&~ (A \vee C)
\]

\subsection{Идемпотентность}
\[
    A \vee A = A
\]

\subsection{Закон де Моргана}
\[
    \neg (A \vee B) = \neg A ~\&~ \neg B
\]

\subsection{Поглощение}
\[
    A \vee (A ~\&~ B) = A
\]

\subsection{Расщепление}
\[
    A \equiv (A \vee B) ~\&~ (A \vee \neg B)
\]

\subsection{Снятие двойного отрицания}
\[
    \neg \neg A \equiv A
\]



\section{Обобщённые законы}
\subsection{Закон де Моргана}
\[
    \neg (A_1 \vee \dots \vee A_n) = \neg A_1 ~\&~ \dots ~\&~ \neg A_n
\]

\subsection{Дистрибутивность}
\begin{gather*}
    (A_1 ~\&~ \dots ~\&~ A_k) \vee (B_1 ~\&~ \dots ~\&~ B_m) \equiv \\
    \equiv (A_1 \vee B_1) ~\&~ (A_1 \vee B_2) ~\&~ \dots ~\&~ (A_1 \vee B_m) ~\&~ (A_2 \vee B_1) ~\&~ \dots ~\&~ (A_k \vee B_m)
\end{gather*}



\section{ДНФ и КНФ}
Дизъюнктивные и конъюнктивные нормальные формы формул

\subsection{Определение}
Элементарная конъюнкция --- конюънкция переменных или их отрицание.

\subsubsection{Пример}
\[
    x, \neg y, x ~\&~ y ~\&~ \neg z
\]


\subsection{Определение}
Элементарная дизъюнкция --- дизъюнкция переменных или их отрицание.

\subsubsection{Пример}
\[
    x, \neg x \vee y \vee z
\]


\subsection{Определение}
$F$ в ДНФ, если $F$ --- дизъюнкция элементарных конъюнкций (возможно, одного)


\subsection{Определение}
$F$ в КНФ, если $F$ --- конъюнкция элементарных дизъюнкций (возможно, одного)

\subsubsection{Пример (для обоих определений)}
\[
    x, \neg y, x ~\&~ y, \neg x \vee z \text{ --- ДНФ, КНФ}
\]

\subsubsection{Пример (для обоих определений)}
\[
    (x ~\&~ \neg y) \vee (x ~\&~ y ~\&~ \neg z) \text{ --- ДНФ}
\]

\subsubsection{Пример (для обоих определений)}
\[
    (\neg x \vee y) ~\&~ (z \vee x) \text{ --- КНФ}
\]



\subsection{Алгоритм ДНФ (КНФ)}
\begin{enumerate}
    \item $\supset$, $\sim$, \dots $\rightarrow$ $\{\&, \vee, \neg\}$
    \item $\neg$ --- тесные (перед переменной, а не перед скобкой) --- закон Де Моргана
    \item Обобщённый закон дистрибутивности
\end{enumerate}


\subsection{Практика}
\subsubsection{Пример}
\begin{gather*}
    \neg(x \supset \neg y) \sim (\neg y \supset \neg z) \equiv
    \neg(\neg x \vee \neg y) \sim (y \vee \neg z) \equiv
    (x ~\&~ y) \sim (y \vee \neg z) \equiv \\
    (x ~\&~ y ~\&~ (y \vee \neg z)) \vee (\neg (x ~\&~ y) ~\&~ \neg (y \vee \neg z)) \equiv
    (x ~\&~ y) \vee ((\neg x \vee \neg y) ~\&~ \neg y ~\&~ z) \equiv
    (x ~\&~ y) \vee (\neg y ~\&~ z)_{--- ДНФ} \equiv \\
    (x \vee \neg y) ~\&~ (x \vee z) ~\&~ (y \vee \neg y) ~\&~ (y \vee z) \equiv
    (x \vee \neg y) ~\&~ (x \vee z) ~\&~ (y \vee z) 
\end{gather*}
\includegraphics[scale=0.3]{images/part08_01.png}


\section{Совершенные ДНФ и КНФ (СДНФ, СКНФ)}
\[
    F <X_1, \dots, X_n> \text{ в СДНФ (СКФН), если}:
\]
\begin{enumerate}
    \item $F$ в ДНФ (КНФ)
    \item $\forall$ элементов \&-ия ($\vee$-ия) упорядочена по всему СП: ~~~~~ $i: X_i$ или $\neg x_i$
    \item $\forall$ элементов \&-ия ($\vee$-ия) ед.
\end{enumerate}

\subsection{Теорема}
\[
    \forall F ~\cancel{\equiv}~ \text{Т.-Л. (Т-И.)} \rightarrow \text{СДНФ (СКНФ). Единств. с точн. до наоборот}
\]

\subsection{Алгоритм}
\begin{enumerate}
    \item $F$ в ДНФ (КНФ)
    \item Закон расщепления и коммутативность
    \item Удаляем все одиночные элементы ~\&~ ($\vee$), кроме одной
\end{enumerate}


\subsection{Пример}
\includegraphics[scale=0.3]{images/part08_02.png}



\paragraph{2024.11.26 (3, 4)}
\section{???}
\subsection{Пример | РГР-5(б)}
\begin{gather*}
    \neg(x \subset y) \sim \neg(z \& \neg x) \equiv
    \neg(\neg x \vee y) \sim \neg(z \& \neg x) \equiv
    (x \& \neg y) \sim (\neg z \vee x) \equiv \\
    \equiv (x \& \neg y (\neg z \vee x)) \vee (\neg(x \& \neg y) \& \neg(\neg z \vee x)) \equiv
    (x \& \neg y) \vee ((\neg x \vee y) \& z \& \neg x) \equiv
    (x \& \neg y) \vee (z \& \neg x) \equiv \\
    \equiv (x \& \neg y \& z) \vee (x \& \neg y \& \neg z) \vee (\neg x \& y \& z) \vee (\neg x \& \neg y \& z) \equiv \\
    \equiv (x \vee z) \& (x \vee \neg x) \& (\neg y \vee z) \& (\neg y \vee \neg x) \equiv
    (x \vee z) \& (\neg y \vee z) \& (\neg y \vee \neg x) \equiv
    (x \vee y \vee z) \& (x \vee \neg y \vee z) \& (x ) ...
\end{gather*}




\part{Булевы функции}

\section{Введение}

\subsection{Определение}
\begin{gather*}
    f(x_1, \dots, x_n): \{0, 1\}^n \rightarrow \{0, 1\} \\
    \text{И} \leftrightarrow 0, \text{Л} \leftrightarrow 1 \\
    \neg x \rightarrow \overline{x}, F \equiv G \rightarrow f = g, + \text{ --- исключающее 'или' } \rightarrow +\mod 2 
\end{gather*}

\subsubsection{Пример | РГР-5(в)}
\[
    (\overline{x \subset y}) \sim (\overline{z \& \overline{x}})
\]
\begin{tabular}{c|c|c|c|c|c|c|c}
    $x$ & $y$ & $z$ & $x \subset y$ & $\overline{x \subset y}$ & $\overline{x}$ & $z \& \overline{x}$ & $(\overline{x \subset y}) \sim (\overline{z \& \overline{x}})$ \\
    \hline
    $1$ & $1$ & $1$ & $1$           & $0$                      & $0$            & $0$                 & $0$ \\
    $1$ & $1$ & $0$ & $1$           & $0$                      & $0$            & $0$                 & $0$ \\
    $1$ & $0$ & $1$ & $0$           & $1$                      & $0$            & $0$                 & $0$ \\
    $1$ & $0$ & $0$ & $0$           & $1$                      & $0$            & $0$                 & $0$ \\
    $0$ & $1$ & $1$ & $1$           & $0$                      & $1$            & $0$                 & $0$ \\
    $0$ & $1$ & $0$ & $1$           & $0$                      & $1$            & $0$                 & $0$ \\
    $0$ & $0$ & $1$ & $1$           & $0$                      & $1$            & $0$                 & $0$ \\
    $0$ & $0$ & $0$ & $1$           & $0$                      & $1$            & $0$                 & $0$ \\
\end{tabular}

\subsection{Теорема | Представление функции в виде формулы в СДНФ}
\[
    \forall f \cancel{\equiv} 0 \text{ м. предст. } F \text{ в СДНФ } \langle x_1, \dots, x_n \rangle. \text{Единств.} \\
\]
\begin{gather*}
    x^1 = x; x^0 = \overline{x} \\
    \langle t \rangle = \langle t_1, \dots, t_n \rangle. \text{Эл. \&-ия} x_1^{t_1} \& x_2^{t_2} \& \dots \& x_n^{t_n} \text{ --- ассоциир. с} \langle t \rangle \\
    \langle 1, 0, 1 \rangle \text{ --- } x_1^1 \& x_2^0 \& x_3^1 = x_1 \& \overline{x_2} \& x_3
\end{gather*}

\subsubsection{Алгоритм $f \to F$ в СДНФ}
\begin{enumerate}
    \item $\forall \langle t \rangle = \langle t_1, \dots, t_n \rangle: f(t_1, \dots, t_n) = 1$.
    \item $\forall \langle t \rangle: f(t) = 1 \leftrightarrow x_1^{t_1} \& x_2^{t_2} \& \dots \& x_n^{t_n}$
    \item $F = \vee ...$ 
\end{enumerate}

\subsubsection{Пример | РГР-5(г)}
\begin{tabular}{c|c|c|c|c}
    $x$ & $y$ & $z$ & $f$   & СДНФ                 \\
    \hline
    $1$ & $1$ & $1$ & $0$   &                      \\  
    $1$ & $1$ & $0$ & $0$   &                      \\
    $1$ & $0$ & $1$ & $1^*$ & $x^1 \& y^0 \& z^1$  \\
    $1$ & $0$ & $0$ & $1^*$ & $x^1 \& y^0 \& z^0$  \\
    $0$ & $1$ & $1$ & $1^*$ & $x^0 \& y^1 \& z^1$  \\
    $0$ & $1$ & $0$ & $0$   &                      \\
    $0$ & $0$ & $1$ & $1^*$ & $x^0 \& y^0 \& z^1$  \\
    $0$ & $0$ & $0$ & $0$   &    
\end{tabular}


\subsection{Теорема | о предет. $f$ в СДНФ}
\begin{gather*}
    \forall f \cancel{\equiv} 1 \text{ м. предет. } F \text{ в СДНФ } \langle x_1, \dots, x_n \rangle. \text{ Единств.} \\
    \langle t \rangle = \langle t_1, \dots, t_n \rangle ~~~ \text{эл. } \vee \text{-ия} x_1^{1 - t_1} \vee x_2^{1 - t_2} \vee \dots \vee x_n^{1 - t_n} \text{ --- ассоц. с } \langle t \rangle \\
    \langle 1, 0, 1 \rangle \text{ --- } x_1^{1 - 1} \vee x_2^{1 - 0} \vee x_3^{1 - 1} = x_1^0 \vee x_2^1 \vee x_3^0 = \overline{x_1} \vee x_2 \vee \overline{x_3}
\end{gather*}

\subsubsection{Алгоритм $f \to$ в СКНФ}
\begin{enumerate}
    \item ...
\end{enumerate}

\begin{tabular}{c|c|c|c|c}
    $x$ & $y$ & $z$ & $f$   & СКНФ                    \\
    \hline
    $1$ & $1$ & $1$ & $0$   & $x^0 \vee y^0 \vee z^0$ \\  
    $1$ & $1$ & $0$ & $0$   & $x^0 \vee y^0 \vee z^1$ \\
    $1$ & $0$ & $1$ & $1^*$ &                         \\
    $1$ & $0$ & $0$ & $1^*$ &                         \\
    $0$ & $1$ & $1$ & $1^*$ &                         \\
    $0$ & $1$ & $0$ & $0$   & $x^1 \vee y^0 \vee z^1$ \\
    $0$ & $0$ & $1$ & $1^*$ &                         \\
    $0$ & $0$ & $0$ & $0$   & $x^1 \vee y^1 \vee z^1$
\end{tabular}



\section{Полные системы булевых функций}

\subsection{Определение}
    $\{f_1, \dots, f_m\}$ --- полная, если для любой булевой функции можно выразить через \\
    $f_1, \dots, f_m$ с помощью суперпозиций. 
\begin{gather*}
    f_1(x_1, \dots, x_{k_1}), \dots, f_m(x_1, \dots, x_{k_m}) \\
    f_z(x_1, \dots, x_{i - 1}, x_i, x_{i + 1}, \dots, x_{k_t}) \\
\end{gather*}

\subsubsection{Достаточное условие полноты}
\begin{gather*}
    \text{Пусть } \{f_1, \dots, f_m\} \text{ --- полная.} \\
    \text{Тогда } \{g_1, \dots, g_k\} \text{ --- полная.} \\
    \text{Если  } \forall f_1, \dots, f_m \text{ можно выразить через } \{g_1, \dots, g_k\} \text{с помощью суперпозиции}. \\
    \{\&, \vee, \neg\} \text{ --- полная} \\
    f_1 \text{ --- } \&, f_2 \text{ --- } \vee, f_3 \text{ --- } \neg \\
\end{gather*}
\paragraph{Доказать полноту $\{\vee, \neg\}$}
\begin{gather*}
    g_1 \text{ --- } \vee, g_2 \text{ --- } \neg \\
    f_2(x, y) \equiv g_1(x, y) \\
    f_3(x) \equiv g_2(x) \\
    x \& y \equiv \neg(\neg x \vee \neg y) \\
    f_1(x, y) = g_2(g_1(g_2(x), g_2(y))) \\
\end{gather*}
\paragraph{Пункт 2}
\begin{gather*}
    \{\subset, 0\} ~~~ g_1 \text{ --- } \subset, g_2 \text{ --- } 0. \\
    \neg x = x \subset 0 \\
    x \vee y = (x \subset 0) \subset y ~~~ (A \subset B \equiv \neg A \vee B \Rightarrow A \vee B \equiv \neg A \subset B) \\
    x \& y = ~~~ (A \subset B \equiv \neg(A \& \neg B) \Rightarrow A \& \neg B \equiv \neg(A \subset B), A \& B \equiv \neg(A \subset \neg B))
\end{gather*}
\paragraph{Пункт 3}
Функция Шеффера
\[
    x \big| y \equiv \neg x \vee \neg y
\]

Функция Вебба
\[
    x \circ y \equiv \neg x \& \neg y
\]

\begin{tabular}{c|c|c|c}
    $x$ & $y$ & $x \big| y$ & $x \circ y$ \\
    \hline
    $1$ & $1$ & $0$ & $0$ \\
    $1$ & $0$ & $1$ & $0$ \\
    $0$ & $1$ & $1$ & $0$ \\
    $0$ & $0$ & $1$ & $1$ 
\end{tabular}


\paragraph{2024.11.26 (4)}
\section{Многочлены Жегалкина}

\subsection{Определение}
\begin{gather*}
    \{\&, +, 1\} \text{ --- полная.} \\
    x \& y \\
    \neg x = x + 1 \\
    x \vee y = \neg(\neg x \& \neg y) = \left( (x + 1) \& (y + 1) \right) + 1 \\
    \& \to \cdot \Rightarrow \{+, \cdot, 1\}
\end{gather*}

\subsection{Свойства}

\subsubsection{Коммутативность}
\[
    x + y = y + x; x \cdot y = y \cdot x
\]

\subsubsection{Ассоциативность}
\[
    x + (y + z) = (x + y) + z; x \cdot (y \cdot z) = (x \cdot y) \cdot x
\]

\subsubsection{...}
\[
    x \cdot (y + z) = x \cdot y + x \cdot z
\]

\subsubsection{...}
\[
    x + x = 0; x \cdot x = x
\]

\begin{tabular}{c|c|c|c|c}
    $x$ & $y$ & $x \& y$ & $x + y$ & $1$ \\
    \hline
    $1$ & $1$ & $1$      & $0$     & $1$ \\
    $1$ & $0$ & $0$      & $1$     & $1$ \\
    $0$ & $1$ & $0$      & $1$     & $1$ \\
    $0$ & $0$ & $0$      & $0$     & $1$
\end{tabular}


\subsection{Определение | Многочлен Жегалкина}
\begin{gather*}
    \Sigma x_{i_1} \cdot x_{i_2} \cdot \dots \cdot x_{i_k} + a_i, \\
    \text{где } a_i = const \in \{0, 1\} \\
    x \vee y = (x + 1) \cdot (y + 1) + 1 =
    xy + x + y + 1 + 1 = xy + x + y \\
    x \subset y = \neg x \vee \neg y = 
    \neg(x \& \neg y) =
    x \cdot (y + 1) + 1 = xy + x + 1 \\
    x \sim y = \neg(x + y) = x + y + 1
\end{gather*}
\begin{gather*}
    x \& y = x \cdot y \\
    \neg x = x + 1 \\
    x + y \\
    x \vee y = xy + x + y \\
    x \subset y = xy + x + 1 \\
    x \sin y = x + y + 1
\end{gather*}

\subsubsection{Пример | РГР-5(д)}
\paragraph{Способ 1}
\[
    F = \neg(x \subset y) \sim \neg(z \neg x) \\
\]
\begin{enumerate}
    \item $\neg(x \subset y) = xy + x + 1 + 1 = xy + x$
    \item $\neg(z \& \neg x) = z(x + 1) + 1 = xz + z + 1$
    \item $F = \neg(x \subset y) \sim \neg(z \neg x) = xy + x + xz + z + 1 + 1 = xy + x + z$
\end{enumerate}
\begin{tabular}{c|c|c|c}
    $x$ & $y$ & $z$ & $f$ \\
    \hline
    $1$ & $1$ & $1$ & $0$ \\
    $1$ & $1$ & $0$ & $0$ \\
    $1$ & $0$ & $1$ & $1$ \\
    $1$ & $0$ & $0$ & $1$ \\
    $0$ & $1$ & $1$ & $1$ \\
    $0$ & $1$ & $1$ & $0$ \\
    $0$ & $0$ & $1$ & $1$ \\
    $0$ & $0$ & $0$ & $0$
\end{tabular}
\begin{gather*}
    f(1, 1, 0) = 0 ~~~ xy + yz + x + z \big|_{\langle 1, 1, 0 \rangle} = 1 \cdot 1 + 1 \cdot 0 + 1 + 0 = 0 \\
    f(1, 0, 1) = 1 ~~~ xy + xz + x + z \big|_{\langle 1, 0, 1 \rangle} = 1 \cdot 0 + 1 \cdot 1 + 1 + 1 = 1
\end{gather*}
\paragraph{Способ 2}
$F$ --- в СДНФ ~~~ $\vee \leftrightarrow +$
\begin{gather*}
    F = x \overline{y} z + x \overline{y} \overline{z} + \overline{x} y z + \overline{x} \overline{y} z = \\
    = x(y + 1)z + x(y + 1)(z + 1) + (x + 1)yz + (x + 1)(y + 1)z = \\
    = xyz + xz + xyz + xy + xz + x + xyz + yz + xyz + xz + yz + x = \\
    = xy + x + xz + z
\end{gather*}



\section{Критерий Поста}

\subsection{Определение}
Класс функций называется функционально-замкнутым (ФЗК), если любые суперпозиции функций этого класса остаются в этом классе.

\subsubsection{Примеры}
\begin{enumerate}
    \item $K$ --- все булевые функции --- ФЗК
    \item $P_1$ --- функции одной переменной --- ФЗК \\
    $f(x), g(x) \in P_1; f(g(x)) \in P_1$
    \item $P_2$ --- функции двух переменных --- не ФЗК \\
    $f_1(x, y), f_2(x, y) \in P_2; f_1(x, f_2(y, z)) \notin P_2$
    \item $T_{\overline{1}}$ --- $f(1, 1, \dots, 1) = 0$ \\
    $+ \in T_{\overline{1}}$ ~~~ $1 + (1 + 1) = 1 \notin T_{\overline{1}}$
\end{enumerate}


\subsection{5 основных ФЗК}
\begin{enumerate}
    \item $T_0$ --- сохраняет $0: f(0, 0, \dots, 0) = 0$ \\
    $\& \in T_0: 0 \& 0 = 0$ \\
    $\subset \notin T_0: 0 \subset 0 = 1$
    \item $T_1$ --- сохраняет $1: f(1, 1, \dots, 1) = 1$ \\
    $\& \in T_1: 1 \& 1 = 1$ \\
    $\subset \notin T_1: 1 \subset 1 = 1$
    \item $S$ --- самодвойственная \\
    $f^*(x_1, \dots, x_n) = \neg f(\neg x_1, \dots, \neg x_n)$ \\
    $f$ --- самодвойственна, если $f = f^*$ \\
    $(x \& y)^* = \neg(\neg x \& \neg y) = x \vee y                          ~~~~~ \&,   \vee \notin S$ \\
    $(x \sim y)^* = \neg(\neg x \sim \neg y) = x + 1 + y + 1 + 1 + 1 = x + y ~~~~~ \sim, +    \notin S$ \\
    $(1)^* = \neg 1 = 0                                                      ~~~~~ 1,    0    \notin S$ \\
    $(\neg x)^* = \neg(\neg(\neg x)) = \neg x                                ~~~~~ \neg       \in    S$ \\
    \item $L$ --- линейные функции \\
    $f \in L$, если $f$ --- линейный множитель Жегалкина. \\
    $f(x_1, \dots, x_n) = a_0 + a_1 x + \cdot + a_n x_n$ \\
    $\vee \notin L: x \vee y = xy + x + y \notin L$ \\
    $\sim \in    L: x \sim y = x + y + 1  \in    L$ \\
    \item $M$ --- монотонные функции \\
    $\langle t \rangle = \langle t_1, \dots, t_n \rangle; \langle s \rangle = \langle s_1, \dots, s_n \rangle$ \\
    $\text{Частичный порядок: } \langle t \rangle \preccurlyeq \langle s \rangle \Leftrightarrow t_i \le s_i, \forall i \in \overline{1, n}$ \\
    $\langle 0, 1, 0 \rangle \preccurlyeq \langle 1, 1, 0 \rangle: \langle 0, 1, 0 \rangle \text{ и } \langle 1, 0, 0 \rangle \text{ --- не сравнимы}$ \\
    $ f(x_1, \dots, x_n) \in M, \text{если } \forall \langle t \rangle \preccurlyeq \langle s \rangle \Rightarrow f(t) \le f(s)$
\end{enumerate}
\begin{tabular}{c|c|c|c}
    $x$ & $y$ & $x \& y$ & $x \sim y$ \\
    \hline
    $1$ & $1$ & $1$      & $1$        \\
    $1$ & $0$ & $0$      & $0$        \\
    $0$ & $1$ & $0$      & $0$        \\
    $0$ & $0$ & $0$      & $1$ 
\end{tabular}


\subsection{Теорема | Критерий Поста}
Для того, чтобы $\{f_1, \dots, f_m\}$ --- полная, необходимо и достаточно, 
чтобы для $\forall T_0, T_1, S, L, M \exists f_i \in \{f\}$, не принадлежащая этому классу.

\subsubsection{Пример}
\[
    \underbrace{\neg x}_{\notin T_0} \neq \underbrace{\&, \vee}_{\in T_0}
\]

\subsubsection{Пример}
\[
    \{\&, \vee, \neg\} \text{ --- полная?}
\]
\begin{tabular}{c|c|c|c|c|c}
    ~         & $T_0$ & $T_1$ & $S$ & $L$ & $M$ \\
    $\&$      &   +   &   +   &  -  &  -  &  +  \\
    $\vee$    &   +   &   +   &  -  &  -  &  +  \\
    $\neg$    &   -   &   -   &  +  &  +  &  -  \\
    $\subset$ &       &       &     &     &     \\
    $\sim$    &       &       &     &     &     \\
    $+$       &   +   &   -   &  -  &  +  &  -  \\
    $1$       &   -   &   +   &  -  &  +  &  +  \\
    $0$       &       &       &     &     &     \\
    $|$       &       &       &     &     &     \\
    $\circ$   &       &       &     &     &     
\end{tabular}
\begin{tabular}{c|c|c|c|c}
    $x$ & $y$ & $x \& y$ & $x \vee y$ & $\neg x$ \\
    $1$ & $1$ & $1$      & $1$        & $0$      \\
    $1$ & $0$ & $0$      & $1$        & $0$      \\
    $0$ & $1$ & $0$      & $1$        & $1$      \\
    $0$ & $0$ & $0$      & $0$        & $1$      
\end{tabular}
\begin{gather*}
    1) \&: \\ 
    0 \& 0 = 0 \in T_0 \\
    1 \& 1 = 1 \in T_1 \\
    (x \& y)^* = \dots = x \vee y \notin S \\
    x \& y = xy \notin L \\
    2) \vee: \\
    0 \vee 0 = 0 \in T_0 \\
    1 \vee 1 = 1 \in T_1 \\
    (x \vee y)^* = x \& y \notin S \\
    x \vee y = xy + x + y \notin L \\
    3) \neg: \\
    \neg 0 = 1 \notin T_0 \\
    \neg 1 = 0 \notin T_1 \\
    (\neg x)^* = \neg x \in S \\
    \neg x = x + 1 \in L
\end{gather*}
\begin{gather*}
    \neg \notin T_0, \neg \notin T_1, \& \notin S, \vee \notin L, \neg \notin M. \\
    \{\vee, \neg\} \text{ --- полная}. \\
    \neg \notin T_0, T_1, M; \vee \notin S, L. \\
    \{\&, \vee\} \text{ --- неполная}: \&, \vee \in T_0 \\
    \{\&, +, 1\} \text{ --- полная}. \\
    \neg \notin T_0, + \notin T_1, + \notin S, \& \notin S, + \notin M. \\
    +: \\
    0 + 0 = 0 \in T_0 \\
    1 + 1 = 0 \notin T_1 \\
    (x + y)^* = x \sim y \notin S \\
    x + y \in L \\
    \langle 0, 1 \rangle \preccurlyeq \langle 1, 1 \rangle, \text{ но } f(0, 1) > f(1, 1)
    \{+, \neg, 1\} \text{ нет}; L \\
    \{\vee, 1\} \text{ --- нет}; T_1
\end{gather*}



\newpage
\section{Независимые системы булевых функций}
\begin{gather*}
    T_0: \neg \notin T_0 ~~~ \underbrace{\neg x}_{\notin T_0} \neq \dots \&, \vee, +v \dots \\
    \{\&, \vee, \neg\} \text{ --- полная} \\
    x \& y \equiv \neg(\neg x \vee \neg y) \Rightarrow \text{ система зависима} \\
    \{\vee, \neg\} \text{ --- полная, независимая?}
\end{gather*}

\subsection{Определение}
$\{f_1, \dots, f_m\}$ --- независимая, если никакая $f$, 



\newpage
\part{Логика предикатов}

\section{Введение}

\subsection{Определение | Предикат}
Предикат --- $n$-местная функция $P^{(n)}_{x_1, \dots, x_n}: M^n = \{$И, Л$\}$

\section{Кванторы}

\subsection{Квантор общности}
$(\forall x)$ --- для любого $x$
\begin{itemize}
    \item $(\forall x) P(x) =$ И $\Leftrightarrow P(x) =$ И для всех $x \in M$.
    \item $(\forall x) P(x) =$ Л, если $P(x) =$ Л для некоторого $x \in M$.
\end{itemize}
$x$ --- связ. перемен.


\subsection{Квантор существования}
$(\exists x)$ --- существование $x$
\begin{itemize}
    \item $(\exists x) P(x) =$ И, если $P(x) =$ И для некоторого $x \in M$.
    \item $(\exists x) P(x) =$ Л $\Leftrightarrow P(x) =$ Л для всех $x \in M$.
\end{itemize}
$x$ --- связ.




\end{document}